%% Generated by Sphinx.
\def\sphinxdocclass{report}
\documentclass[letterpaper,10pt,english]{sphinxmanual}
\ifdefined\pdfpxdimen
   \let\sphinxpxdimen\pdfpxdimen\else\newdimen\sphinxpxdimen
\fi \sphinxpxdimen=.75bp\relax
\ifdefined\pdfimageresolution
    \pdfimageresolution= \numexpr \dimexpr1in\relax/\sphinxpxdimen\relax
\fi
%% let collapsible pdf bookmarks panel have high depth per default
\PassOptionsToPackage{bookmarksdepth=5}{hyperref}

\PassOptionsToPackage{booktabs}{sphinx}
\PassOptionsToPackage{colorrows}{sphinx}

\PassOptionsToPackage{warn}{textcomp}
\usepackage[utf8]{inputenc}
\ifdefined\DeclareUnicodeCharacter
% support both utf8 and utf8x syntaxes
  \ifdefined\DeclareUnicodeCharacterAsOptional
    \def\sphinxDUC#1{\DeclareUnicodeCharacter{"#1}}
  \else
    \let\sphinxDUC\DeclareUnicodeCharacter
  \fi
  \sphinxDUC{00A0}{\nobreakspace}
  \sphinxDUC{2500}{\sphinxunichar{2500}}
  \sphinxDUC{2502}{\sphinxunichar{2502}}
  \sphinxDUC{2514}{\sphinxunichar{2514}}
  \sphinxDUC{251C}{\sphinxunichar{251C}}
  \sphinxDUC{2572}{\textbackslash}
\fi
\usepackage{cmap}
\usepackage[T1]{fontenc}
\usepackage{amsmath,amssymb,amstext}
\usepackage{babel}



\usepackage{tgtermes}
\usepackage{tgheros}
\renewcommand{\ttdefault}{txtt}



\usepackage[Bjarne]{fncychap}
\usepackage{sphinx}

\fvset{fontsize=auto}
\usepackage{geometry}


% Include hyperref last.
\usepackage{hyperref}
% Fix anchor placement for figures with captions.
\usepackage{hypcap}% it must be loaded after hyperref.
% Set up styles of URL: it should be placed after hyperref.
\urlstyle{same}

\addto\captionsenglish{\renewcommand{\contentsname}{Contents:}}

\usepackage{sphinxmessages}
\setcounter{tocdepth}{1}



\title{xmmPipeline}
\date{Feb 13, 2023}
\release{}
\author{Suyog}
\newcommand{\sphinxlogo}{\vbox{}}
\renewcommand{\releasename}{}
\makeindex
\begin{document}

\ifdefined\shorthandoff
  \ifnum\catcode`\=\string=\active\shorthandoff{=}\fi
  \ifnum\catcode`\"=\active\shorthandoff{"}\fi
\fi

\pagestyle{empty}
\sphinxmaketitle
\pagestyle{plain}
\sphinxtableofcontents
\pagestyle{normal}
\phantomsection\label{\detokenize{index::doc}}


\sphinxstepscope


\chapter{xmmPipeline}
\label{\detokenize{modules:xmmpipeline}}\label{\detokenize{modules::doc}}
\sphinxstepscope


\section{xmmPipeline package}
\label{\detokenize{xmmPipeline:xmmpipeline-package}}\label{\detokenize{xmmPipeline::doc}}

\subsection{Submodules}
\label{\detokenize{xmmPipeline:submodules}}

\subsection{xmmPipeline.convert module}
\label{\detokenize{xmmPipeline:module-xmmPipeline.convert}}\label{\detokenize{xmmPipeline:xmmpipeline-convert-module}}\index{module@\spxentry{module}!xmmPipeline.convert@\spxentry{xmmPipeline.convert}}\index{xmmPipeline.convert@\spxentry{xmmPipeline.convert}!module@\spxentry{module}}
\sphinxAtStartPar
For conversion of the RA and DEC.


\subsection{xmmPipeline.epicObj module}
\label{\detokenize{xmmPipeline:module-xmmPipeline.epicObj}}\label{\detokenize{xmmPipeline:xmmpipeline-epicobj-module}}\index{module@\spxentry{module}!xmmPipeline.epicObj@\spxentry{xmmPipeline.epicObj}}\index{xmmPipeline.epicObj@\spxentry{xmmPipeline.epicObj}!module@\spxentry{module}}
\sphinxAtStartPar
10th May 2021:
—
Am making several changes to the code procedure arrangements.
See the GitHub repo, the Notes on Google Docs and the documentation for more information.
For previous docstring comments, see earlier code files, namely \sphinxtitleref{xmmPipeline.py}

\sphinxAtStartPar
18th and 19th May 2021:
—
Adding options for Pile\sphinxhyphen{}up obsIDs.

\sphinxAtStartPar
21st May 2021:
—
Editing the \sphinxtitleref{removeFlareBackground} function.
A seperate function \sphinxtitleref{extract\_flareGTI} will also be added.

\sphinxAtStartPar
9th\textasciitilde{}13th June 2021:
—
Starting to write the functions for obtaining individual background circles.

\sphinxAtStartPar
20th Jan 2023:
—
Making the code compatible with Mac.
If the path names are Windows like, containing white\sphinxhyphen{}spaces, then we need to have
additional single quotes around them before calling the \sphinxtitleref{subprocess.run} command.
\index{epicObj (class in xmmPipeline.epicObj)@\spxentry{epicObj}\spxextra{class in xmmPipeline.epicObj}}

\begin{fulllineitems}
\phantomsection\label{\detokenize{xmmPipeline:xmmPipeline.epicObj.epicObj}}
\pysigstartsignatures
\pysiglinewithargsret{\sphinxbfcode{\sphinxupquote{class\DUrole{w}{  }}}\sphinxcode{\sphinxupquote{xmmPipeline.epicObj.}}\sphinxbfcode{\sphinxupquote{epicObj}}}{\emph{\DUrole{n}{ra}}, \emph{\DUrole{n}{dec}}, \emph{\DUrole{n}{workdir}}, \emph{\DUrole{n}{sas\_dir}}, \emph{\DUrole{n}{headas}}, \emph{\DUrole{n}{sas\_ccfpath}}, \emph{\DUrole{n}{srcCircRadius}\DUrole{o}{=}\DUrole{default_value}{None}}, \emph{\DUrole{n}{bkgCircRadius}\DUrole{o}{=}\DUrole{default_value}{None}}, \emph{\DUrole{n}{dSrcThreshold}\DUrole{o}{=}\DUrole{default_value}{70}}, \emph{\DUrole{n}{bkgCircGap}\DUrole{o}{=}\DUrole{default_value}{2.5}}, \emph{\DUrole{n}{edetectmode}\DUrole{o}{=}\DUrole{default_value}{\textquotesingle{}chain\textquotesingle{}}}, \emph{\DUrole{n}{esp\_nsplinenodes}\DUrole{o}{=}\DUrole{default_value}{14}}, \emph{\DUrole{n}{gti\_indiThreshold}\DUrole{o}{=}\DUrole{default_value}{500}}, \emph{\DUrole{n}{gti\_combThreshold}\DUrole{o}{=}\DUrole{default_value}{None}}, \emph{\DUrole{n}{lcBinSize}\DUrole{o}{=}\DUrole{default_value}{25}}, \emph{\DUrole{n}{binBkglc}\DUrole{o}{=}\DUrole{default_value}{\textquotesingle{}no\textquotesingle{}}}, \emph{\DUrole{n}{saveFig}\DUrole{o}{=}\DUrole{default_value}{True}}, \emph{\DUrole{n}{showFig}\DUrole{o}{=}\DUrole{default_value}{True}}, \emph{\DUrole{n}{ignorePileup}\DUrole{o}{=}\DUrole{default_value}{False}}}{}
\pysigstopsignatures
\sphinxAtStartPar
Bases: \sphinxhref{https://docs.python.org/3/library/functions.html\#object}{\sphinxcode{\sphinxupquote{object}}}
\index{combineEPICdata() (xmmPipeline.epicObj.epicObj method)@\spxentry{combineEPICdata()}\spxextra{xmmPipeline.epicObj.epicObj method}}

\begin{fulllineitems}
\phantomsection\label{\detokenize{xmmPipeline:xmmPipeline.epicObj.epicObj.combineEPICdata}}
\pysigstartsignatures
\pysiglinewithargsret{\sphinxbfcode{\sphinxupquote{combineEPICdata}}}{}{}
\pysigstopsignatures
\sphinxAtStartPar
Combines EPIC PN, MOS1 and MOS2 Event lists data corresponding to the CCD
where the Source is located.
\begin{quote}\begin{description}
\sphinxlineitem{Input}
\sphinxAtStartPar
EPIC concatenated and calibrated Event lists obtained by running
emproc and epproc.

\sphinxlineitem{Output}
\sphinxAtStartPar
Combined PNMOS12.evts and MOS12.evts Event lists.

\end{description}\end{quote}

\end{fulllineitems}

\index{combineGTIs() (xmmPipeline.epicObj.epicObj method)@\spxentry{combineGTIs()}\spxextra{xmmPipeline.epicObj.epicObj method}}

\begin{fulllineitems}
\phantomsection\label{\detokenize{xmmPipeline:xmmPipeline.epicObj.epicObj.combineGTIs}}
\pysigstartsignatures
\pysiglinewithargsret{\sphinxbfcode{\sphinxupquote{combineGTIs}}}{}{}
\pysigstopsignatures
\sphinxAtStartPar
Combines Flare and Instrumental GTIs.
\begin{quote}\begin{description}
\sphinxlineitem{Input}
\sphinxAtStartPar
Flare, EPIC PN, MOS1 and MOS2 GTIs.

\sphinxlineitem{Output}
\sphinxAtStartPar
Combined Flare and Instrumental GTI FITS file.

\end{description}\end{quote}

\end{fulllineitems}

\index{downloadData() (xmmPipeline.epicObj.epicObj method)@\spxentry{downloadData()}\spxextra{xmmPipeline.epicObj.epicObj method}}

\begin{fulllineitems}
\phantomsection\label{\detokenize{xmmPipeline:xmmPipeline.epicObj.epicObj.downloadData}}
\pysigstartsignatures
\pysiglinewithargsret{\sphinxbfcode{\sphinxupquote{downloadData}}}{}{}
\pysigstopsignatures\begin{quote}\begin{description}
\sphinxlineitem{Input}
\sphinxAtStartPar
List of obsIDs for which data has to be downloaded.

\sphinxlineitem{Output}
\sphinxAtStartPar
None.

\end{description}\end{quote}

\sphinxAtStartPar
First it is checked is the the downloading has been done before.
If yes, then if the size of the obsID folder is large, it’s downloading is skipped.

\end{fulllineitems}

\index{extract\_InstrumentalGTIs() (xmmPipeline.epicObj.epicObj method)@\spxentry{extract\_InstrumentalGTIs()}\spxextra{xmmPipeline.epicObj.epicObj method}}

\begin{fulllineitems}
\phantomsection\label{\detokenize{xmmPipeline:xmmPipeline.epicObj.epicObj.extract_InstrumentalGTIs}}
\pysigstartsignatures
\pysiglinewithargsret{\sphinxbfcode{\sphinxupquote{extract\_InstrumentalGTIs}}}{}{}
\pysigstopsignatures
\sphinxAtStartPar
Extracts Instrumental GTIs from EPIC data using ‘fextract’ command.
\begin{quote}\begin{description}
\sphinxlineitem{Input}
\sphinxAtStartPar
EPIC concatenated and calibrated Event lists obtained by running
emproc and epproc.

\sphinxlineitem{Output}
\sphinxAtStartPar
Intrumental GTI FITS files for EPIC PN, MOS1 and MOS2.

\end{description}\end{quote}

\sphinxAtStartPar
‘fextract’ command can’t overwrite the extracted file if it already exists
in the work directory, so they are deleted using ‘rm \sphinxhyphen{}rf’, which doesn’t
raises any error even if no such files are found for deleting during the
first run.

\end{fulllineitems}

\index{extract\_flareGTI() (xmmPipeline.epicObj.epicObj method)@\spxentry{extract\_flareGTI()}\spxextra{xmmPipeline.epicObj.epicObj method}}

\begin{fulllineitems}
\phantomsection\label{\detokenize{xmmPipeline:xmmPipeline.epicObj.epicObj.extract_flareGTI}}
\pysigstartsignatures
\pysiglinewithargsret{\sphinxbfcode{\sphinxupquote{extract\_flareGTI}}}{}{}
\pysigstopsignatures
\sphinxAtStartPar
Runs the XMMSAS commands to extract Flare GTI from PN Event List.
\begin{quote}\begin{description}
\sphinxlineitem{Input}
\sphinxAtStartPar
Unfiltered EPIC PN and MOS1\&2 Events lists.

\sphinxlineitem{Output}
\sphinxAtStartPar
\sphinxtitleref{flareGTI} file for the obsIDs.

\end{description}\end{quote}

\sphinxAtStartPar
This function largely follows the following SAS thread:
\sphinxurl{https://www.cosmos.esa.int/web/xmm-newton/sas-thread-epic-filterbackground}

\sphinxAtStartPar
Only one of either EPIC PN or MOS Flare Background filtering is required
since the Flare affects the whole telescope.

\end{fulllineitems}

\index{extract\_srcBkg\_eventLists() (xmmPipeline.epicObj.epicObj method)@\spxentry{extract\_srcBkg\_eventLists()}\spxextra{xmmPipeline.epicObj.epicObj method}}

\begin{fulllineitems}
\phantomsection\label{\detokenize{xmmPipeline:xmmPipeline.epicObj.epicObj.extract_srcBkg_eventLists}}
\pysigstartsignatures
\pysiglinewithargsret{\sphinxbfcode{\sphinxupquote{extract\_srcBkg\_eventLists}}}{}{}
\pysigstopsignatures
\sphinxAtStartPar
Extracts Source and Background Event list from the combined PNMOS12 data file,
using the Source and Background location parameters given as Input.
These Event lists will then be used for getting source images in JPEG using DS9.

\sphinxAtStartPar
19th May 2021: Added options for Pile\sphinxhyphen{}up obsIDs.
\begin{quote}\begin{description}
\sphinxlineitem{Input}
\sphinxAtStartPar
Combined PNMOS12 Event list, Source and Background X, Y in Sky coords and
rIn, rOut, the radii of the circles.

\sphinxlineitem{Output}
\sphinxAtStartPar
Extracted Source and Background Event lists.

\end{description}\end{quote}

\end{fulllineitems}

\index{findObsIDs() (xmmPipeline.epicObj.epicObj method)@\spxentry{findObsIDs()}\spxextra{xmmPipeline.epicObj.epicObj method}}

\begin{fulllineitems}
\phantomsection\label{\detokenize{xmmPipeline:xmmPipeline.epicObj.epicObj.findObsIDs}}
\pysigstartsignatures
\pysiglinewithargsret{\sphinxbfcode{\sphinxupquote{findObsIDs}}}{}{}
\pysigstopsignatures\begin{quote}\begin{description}
\sphinxlineitem{Input}
\sphinxAtStartPar
Coordinates of the object and path to work directory.

\sphinxlineitem{Output}
\sphinxAtStartPar
List of obsIDs.

\end{description}\end{quote}

\end{fulllineitems}

\index{findSourceCCD() (xmmPipeline.epicObj.epicObj method)@\spxentry{findSourceCCD()}\spxextra{xmmPipeline.epicObj.epicObj method}}

\begin{fulllineitems}
\phantomsection\label{\detokenize{xmmPipeline:xmmPipeline.epicObj.epicObj.findSourceCCD}}
\pysigstartsignatures
\pysiglinewithargsret{\sphinxbfcode{\sphinxupquote{findSourceCCD}}}{}{}
\pysigstopsignatures
\sphinxAtStartPar
Uses ‘ecoordconv’ command of SAS to find which CCD contains the Source
in each of the EPIC’s three cameras.
\begin{quote}\begin{description}
\sphinxlineitem{Input}
\sphinxAtStartPar
EPIC concatenated and calibrated Event lists obtained by running
emproc and epproc.

\sphinxlineitem{Output}
\sphinxAtStartPar
None.

\end{description}\end{quote}

\sphinxAtStartPar
self.sourceCCDs is updated as a dictionary of dictionaries containing the
Source CCD number for PN, MOS1 and MOS2 cameras for all the obsIDs.

\end{fulllineitems}

\index{find\_otherSources() (xmmPipeline.epicObj.epicObj method)@\spxentry{find\_otherSources()}\spxextra{xmmPipeline.epicObj.epicObj method}}

\begin{fulllineitems}
\phantomsection\label{\detokenize{xmmPipeline:xmmPipeline.epicObj.epicObj.find_otherSources}}
\pysigstartsignatures
\pysiglinewithargsret{\sphinxbfcode{\sphinxupquote{find\_otherSources}}}{}{}
\pysigstopsignatures
\sphinxAtStartPar
Function to find all the other Sources, in addition to the main Source, lying within the region of interest.
See: \sphinxurl{https://www.cosmos.esa.int/web/xmm-newton/sas-thread-src-find-stepbystep}
\begin{quote}\begin{description}
\sphinxlineitem{Input}
\sphinxAtStartPar
PNMOS12.evts and PN\_CCD\#\#.evts and corresponding extracted images in the Energy band 0.3\sphinxhyphen{}10 KeV
using evselect.

\sphinxlineitem{Output}
\sphinxAtStartPar
Exposer maps, masks, the csv file containing the x and y coordinates of the other Sources,
apart from the region file from DS9.

\end{description}\end{quote}

\sphinxAtStartPar
Note that other sources may be found within any of the regions using different input files,
for instance ‘PN\_image\_full.fits’ obtained from ‘PNclean.ds’ etc.

\sphinxAtStartPar
However, only the sources within the PNMOS12 overlap region or the PN region, in case of small mode,
are of interest. Since, the PN region necessarily contains the overlap region within it regardless of
the mode of observation, this should have been sufficient to get other sources. However, the MOS data adds
some of its own sources and ‘edetect\_chain’ is easier to use and can be used with multiple imagesets,
so both ‘PNMOS12’ and ‘PN\_CCD\#\#’ are being used to get the sources.

\sphinxAtStartPar
Further, right now the SAS command ‘edetect\_chain’, which is pretty slow when number of sources is large,
is been used. Individual tasks within ‘edetect\_chain’ can be run separately to fasten this up, although
that would require changes in the createOutputImages() function below.

\end{fulllineitems}

\index{find\_otherSources\_indi() (xmmPipeline.epicObj.epicObj method)@\spxentry{find\_otherSources\_indi()}\spxextra{xmmPipeline.epicObj.epicObj method}}

\begin{fulllineitems}
\phantomsection\label{\detokenize{xmmPipeline:xmmPipeline.epicObj.epicObj.find_otherSources_indi}}
\pysigstartsignatures
\pysiglinewithargsret{\sphinxbfcode{\sphinxupquote{find\_otherSources\_indi}}}{\emph{\DUrole{n}{inst}\DUrole{o}{=}\DUrole{default_value}{\textquotesingle{}PN\textquotesingle{}}}}{}
\pysigstopsignatures
\sphinxAtStartPar
Function to find all the other Sources, in addition to the main Source, lying within the region of interest.
See: \sphinxurl{https://www.cosmos.esa.int/web/xmm-newton/sas-thread-src-find-stepbystep}
\begin{quote}\begin{description}
\sphinxlineitem{Parameters}
\sphinxAtStartPar
\sphinxstyleliteralstrong{\sphinxupquote{inst}} (\sphinxhref{https://docs.python.org/3/library/stdtypes.html\#str}{\sphinxstyleliteralemphasis{\sphinxupquote{str}}}) \textendash{} name of the instrument of use.

\sphinxlineitem{Input}
\sphinxAtStartPar
\sphinxtitleref{inst\_CCD\#\#.evts} Event list.

\sphinxlineitem{Raises}
\sphinxAtStartPar
\sphinxstyleliteralstrong{\sphinxupquote{Error}} \textendash{} if \sphinxtitleref{inst} is not in {[}‘PN’, ‘MOS1’, ‘MOS2’{]}

\sphinxlineitem{Output}
\sphinxAtStartPar
Exposer maps, masks, the csv file containing the x and y coordinates of the other Sources,
apart from the region file from DS9.
Also extracts \sphinxtitleref{inst\_image\_CCD\#\#.fits} file which is later used to find srcBkg circs.
Updates the \sphinxcode{\sphinxupquote{otherSourceIndi}} dictionary for each of the obsIDs in \sphinxcode{\sphinxupquote{self.obsIDs}}

\end{description}\end{quote}
\subsubsection*{Notes}

\sphinxAtStartPar
This separate function based on \sphinxtitleref{find\_otherSources} is meant to be used for
finding individual instrument other sources. It uses \sphinxcode{\sphinxupquote{edetect\_chain}} function
from XMMSAS for detecting the Sources.
\sphinxtitleref{inst\_CCD\#\#.evts} is required to extract corresponding images in the Energy band
0.3\sphinxhyphen{}10 KeV using \sphinxcode{\sphinxupquote{evselect}}.
\begin{itemize}
\item {} 
\sphinxAtStartPar
The indi inst files have \sphinxcode{\sphinxupquote{inst}} appended at the end of their names.

\item {} 
\sphinxAtStartPar
\sphinxtitleref{emllist.fits} originally found for the overlap is rewritten with every run
of \sphinxcode{\sphinxupquote{edetect\_chain}}. However, the \sphinxtitleref{emllist.csv} is not b’cuz of the above point
So \sphinxtitleref{emllist.csv} should be used for the overlap cases.
And anyway, another run of the original \sphinxtitleref{find\_otherSources} will update the said file.

\end{itemize}

\begin{sphinxadmonition}{note}{\label{\detokenize{xmmPipeline:id1}}Todo:}
\sphinxAtStartPar
Check is using \sphinxcode{\sphinxupquote{PATTERN in {[}0:12{]}}} in \sphinxcode{\sphinxupquote{evselect expression}} for PN is correct?
For data extraction elsewhere, the expression for PN reads, \sphinxcode{\sphinxupquote{(PATTERN \textless{}= 4)}}
\end{sphinxadmonition}

\end{fulllineitems}

\index{getBackgroundCircles() (xmmPipeline.epicObj.epicObj method)@\spxentry{getBackgroundCircles()}\spxextra{xmmPipeline.epicObj.epicObj method}}

\begin{fulllineitems}
\phantomsection\label{\detokenize{xmmPipeline:xmmPipeline.epicObj.epicObj.getBackgroundCircles}}
\pysigstartsignatures
\pysiglinewithargsret{\sphinxbfcode{\sphinxupquote{getBackgroundCircles}}}{}{}
\pysigstopsignatures
\sphinxAtStartPar
Calls \sphinxtitleref{findOverlap.py} functions to automatically detect the overlap PNMOS12 region and
find 2 background cicles in it.
If the MOS mode is \sphinxtitleref{small}, then the PN image is used for getting the background circles.
\begin{quote}\begin{description}
\sphinxlineitem{Input}
\sphinxAtStartPar
PN\_image.fits, MOS12\_image.fits and PNMOS12\_image.fits to get the overlap from.

\sphinxlineitem{Output}
\sphinxAtStartPar
x, y coords of the background circles alongwith the radii of these circles and the
source circle. Units of radii are pixels or Sky Coord units.
Updates the \sphinxcode{\sphinxupquote{sourceLoc}} and \sphinxcode{\sphinxupquote{backgroundLoc}} dictionaries with the location
parameters.

\end{description}\end{quote}

\end{fulllineitems}

\index{getBackgroundCircles\_indi() (xmmPipeline.epicObj.epicObj method)@\spxentry{getBackgroundCircles\_indi()}\spxextra{xmmPipeline.epicObj.epicObj method}}

\begin{fulllineitems}
\phantomsection\label{\detokenize{xmmPipeline:xmmPipeline.epicObj.epicObj.getBackgroundCircles_indi}}
\pysigstartsignatures
\pysiglinewithargsret{\sphinxbfcode{\sphinxupquote{getBackgroundCircles\_indi}}}{\emph{\DUrole{n}{inst}\DUrole{o}{=}\DUrole{default_value}{\textquotesingle{}PN\textquotesingle{}}}}{}
\pysigstopsignatures
\sphinxAtStartPar
Calls the \sphinxtitleref{backgroundCircles} function in \sphinxtitleref{findOverlap.py} to automatically
find 2 background cicles in the data of the instrument passed.
\begin{quote}\begin{description}
\sphinxlineitem{Parameters}
\sphinxAtStartPar
\sphinxstyleliteralstrong{\sphinxupquote{inst}} (\sphinxhref{https://docs.python.org/3/library/stdtypes.html\#str}{\sphinxstyleliteralemphasis{\sphinxupquote{str}}}) \textendash{} name of the instrument of use.

\sphinxlineitem{Input}
\sphinxAtStartPar
\sphinxtitleref{inst\_CCD\#\#\_image.fits} files containing the extracted image data of the instrument.

\sphinxlineitem{Requires}
\sphinxAtStartPar
Call to \sphinxtitleref{find\_otherSources\_indi} function for the dictionary of
other Sources.

\sphinxlineitem{Raises}
\sphinxAtStartPar
\sphinxstyleliteralstrong{\sphinxupquote{Error}} \textendash{} if \sphinxtitleref{inst} is not in {[}‘PN’, ‘MOS1’, ‘MOS2’{]}

\sphinxlineitem{Output}
\sphinxAtStartPar
x, y coords of the background circles alongwith the radii of these circles.
Units of radii are pixels or Sky Coord units.
Updates the \sphinxcode{\sphinxupquote{backgroundLocIndi}} dictionaries with the background location
parameters.

\end{description}\end{quote}
\subsubsection*{Notes}

\sphinxAtStartPar
\sphinxcode{\sphinxupquote{small\sphinxhyphen{}mode}} MOS checking is invalid here since overlap is not been detected.
The CCD argument in \sphinxcode{\sphinxupquote{findOverlap}} class is named \sphinxcode{\sphinxupquote{pnCCD}}.

\end{fulllineitems}

\index{obtain\_lightCurves() (xmmPipeline.epicObj.epicObj method)@\spxentry{obtain\_lightCurves()}\spxextra{xmmPipeline.epicObj.epicObj method}}

\begin{fulllineitems}
\phantomsection\label{\detokenize{xmmPipeline:xmmPipeline.epicObj.epicObj.obtain_lightCurves}}
\pysigstartsignatures
\pysiglinewithargsret{\sphinxbfcode{\sphinxupquote{obtain\_lightCurves}}}{}{}
\pysigstopsignatures
\sphinxAtStartPar
Obtains the light curves from the Source and Background Event lists.
\begin{quote}\begin{description}
\sphinxlineitem{Input}
\sphinxAtStartPar
Source and Background Event lists.

\sphinxlineitem{Output}
\sphinxAtStartPar
Corresponding light curves.

\end{description}\end{quote}

\sphinxAtStartPar
By default, the background light curve is not binned.

\end{fulllineitems}

\index{readCCDcoordsPickle() (xmmPipeline.epicObj.epicObj method)@\spxentry{readCCDcoordsPickle()}\spxextra{xmmPipeline.epicObj.epicObj method}}

\begin{fulllineitems}
\phantomsection\label{\detokenize{xmmPipeline:xmmPipeline.epicObj.epicObj.readCCDcoordsPickle}}
\pysigstartsignatures
\pysiglinewithargsret{\sphinxbfcode{\sphinxupquote{readCCDcoordsPickle}}}{}{}
\pysigstopsignatures
\sphinxAtStartPar
Reads \sphinxtitleref{ccd\_coords\_info.pickle} file which contains the location parameters.

\end{fulllineitems}

\index{reduceEPICdata() (xmmPipeline.epicObj.epicObj method)@\spxentry{reduceEPICdata()}\spxextra{xmmPipeline.epicObj.epicObj method}}

\begin{fulllineitems}
\phantomsection\label{\detokenize{xmmPipeline:xmmPipeline.epicObj.epicObj.reduceEPICdata}}
\pysigstartsignatures
\pysiglinewithargsret{\sphinxbfcode{\sphinxupquote{reduceEPICdata}}}{}{}
\pysigstopsignatures
\sphinxAtStartPar
Runs the XMMSAS commands on the shell to reduce EPIC data.
\begin{quote}\begin{description}
\sphinxlineitem{Input}
\sphinxAtStartPar
SAS Summary file in the work directory, that is downloaded alongwith the data.

\sphinxlineitem{Output}
\sphinxAtStartPar
EPIC PN and MOS1\&2 calibrated and concatenated Event lists
produced by SAS tasks \sphinxtitleref{epproc} and \sphinxtitleref{emproc}.

\end{description}\end{quote}

\sphinxAtStartPar
These event lists contain Instrumental GTI for each of the CCDs.
See Notes and \sphinxurl{http://xmm-tools.cosmos.esa.int/external/sas/current/doc/epicproc/node17.html}

\sphinxAtStartPar
This function largely follows the following SAS thread:
\sphinxurl{https://www.cosmos.esa.int/web/xmm-newton/sas-thread-epic-reprocessing}

\end{fulllineitems}

\index{removeFlareBackground() (xmmPipeline.epicObj.epicObj method)@\spxentry{removeFlareBackground()}\spxextra{xmmPipeline.epicObj.epicObj method}}

\begin{fulllineitems}
\phantomsection\label{\detokenize{xmmPipeline:xmmPipeline.epicObj.epicObj.removeFlareBackground}}
\pysigstartsignatures
\pysiglinewithargsret{\sphinxbfcode{\sphinxupquote{removeFlareBackground}}}{}{}
\pysigstopsignatures
\sphinxAtStartPar
Runs the XMMSAS commands to remove Flare background from EPIC Event List
and obtained cleaned event lists for each EPIC instrument.
The Flare GTI file extracted previously using \sphinxtitleref{extract\_flareGTI} is used.
Since the Flare Background is same for all the instruments onboard,
this same \sphinxtitleref{flareGTI} file can be used for them all.

\sphinxAtStartPar
Note that these Flare Background filtered Event Lists are only used for
Spectral Analysis, since for Light Curve extraction, \sphinxtitleref{combinedGTI} file
and combined \sphinxtitleref{PNMOS12.evts} EPIC data is used.
\begin{quote}\begin{description}
\sphinxlineitem{Input}
\sphinxAtStartPar
Unfiltered EPIC PN and MOS1\&2 Events lists and \sphinxtitleref{flareGTI} file.

\sphinxlineitem{Output}
\sphinxAtStartPar
Flare Background filtered Event Lists for each EPIC instrument.

\end{description}\end{quote}

\sphinxAtStartPar
This function largely follows the following SAS thread:
\sphinxurl{https://www.cosmos.esa.int/web/xmm-newton/sas-thread-epic-filterbackground}

\end{fulllineitems}

\index{save\_results() (xmmPipeline.epicObj.epicObj method)@\spxentry{save\_results()}\spxextra{xmmPipeline.epicObj.epicObj method}}

\begin{fulllineitems}
\phantomsection\label{\detokenize{xmmPipeline:xmmPipeline.epicObj.epicObj.save_results}}
\pysigstartsignatures
\pysiglinewithargsret{\sphinxbfcode{\sphinxupquote{save\_results}}}{}{}
\pysigstopsignatures\begin{description}
\sphinxlineitem{Does two tasks,}\begin{description}
\sphinxlineitem{One, saves a pickle file at each obsID directory containing}
\sphinxAtStartPar
sourceCCDs, sourceLoc, backgroundLoc and otherSources for it.

\sphinxlineitem{Two, copies the Source, Background Events lists and other output files}
\sphinxAtStartPar
from each obsID directory to a results folder in the main directory.

\end{description}

\end{description}

\end{fulllineitems}

\index{sortObsIDs() (xmmPipeline.epicObj.epicObj method)@\spxentry{sortObsIDs()}\spxextra{xmmPipeline.epicObj.epicObj method}}

\begin{fulllineitems}
\phantomsection\label{\detokenize{xmmPipeline:xmmPipeline.epicObj.epicObj.sortObsIDs}}
\pysigstartsignatures
\pysiglinewithargsret{\sphinxbfcode{\sphinxupquote{sortObsIDs}}}{\emph{\DUrole{n}{objName}\DUrole{o}{=}\DUrole{default_value}{None}}}{}
\pysigstopsignatures
\sphinxAtStartPar
To sort the obsIDs in \sphinxtitleref{obsIDs\_list.dat} by reading it into a pandas dataframe.

\end{fulllineitems}

\index{updateCCDcoordsPickle() (xmmPipeline.epicObj.epicObj method)@\spxentry{updateCCDcoordsPickle()}\spxextra{xmmPipeline.epicObj.epicObj method}}

\begin{fulllineitems}
\phantomsection\label{\detokenize{xmmPipeline:xmmPipeline.epicObj.epicObj.updateCCDcoordsPickle}}
\pysigstartsignatures
\pysiglinewithargsret{\sphinxbfcode{\sphinxupquote{updateCCDcoordsPickle}}}{\emph{\DUrole{n}{inst}\DUrole{o}{=}\DUrole{default_value}{None}}}{}
\pysigstopsignatures
\sphinxAtStartPar
Writes the background location parameter dictionary for individual instruments
\sphinxcode{\sphinxupquote{backgroundLoc\_indi}}  to \sphinxtitleref{ccd\_coords\_info.pickle} file.
If the file is already present in the directory, then the values for the
corresponding keys are updated.
\begin{quote}\begin{description}
\sphinxlineitem{Parameters}
\sphinxAtStartPar
\sphinxstyleliteralstrong{\sphinxupquote{inst}} (\sphinxhref{https://docs.python.org/3/library/stdtypes.html\#str}{\sphinxstyleliteralemphasis{\sphinxupquote{str}}}) \textendash{} name of the instrument of use.

\sphinxlineitem{Input}
\sphinxAtStartPar
\sphinxcode{\sphinxupquote{backgroundLoc\_indi}} dictionary containing the location parameters.

\sphinxlineitem{Output}
\sphinxAtStartPar
Updated \sphinxcode{\sphinxupquote{ccd\_coords\_info.pickle}} file.

\end{description}\end{quote}
\subsubsection*{Notes}

\sphinxAtStartPar
If no CCDcoordsPickle file is available in the \sphinxcode{\sphinxupquote{workdir}}, then the whole
file is created with other values been saved as well.

\end{fulllineitems}

\index{writeCCDcoordsPickle() (xmmPipeline.epicObj.epicObj method)@\spxentry{writeCCDcoordsPickle()}\spxextra{xmmPipeline.epicObj.epicObj method}}

\begin{fulllineitems}
\phantomsection\label{\detokenize{xmmPipeline:xmmPipeline.epicObj.epicObj.writeCCDcoordsPickle}}
\pysigstartsignatures
\pysiglinewithargsret{\sphinxbfcode{\sphinxupquote{writeCCDcoordsPickle}}}{}{}
\pysigstopsignatures
\sphinxAtStartPar
Writes the location parameters to \sphinxtitleref{ccd\_coords\_info.pickle} file.
If the file is already present in the directory, then the values for the
corresponding keys will be updated.
\begin{description}
\sphinxlineitem{21st May 2021: Added loading \sphinxtitleref{ccd\_coords\_info.pickle} file if it is already}
\sphinxAtStartPar
present in the directory.

\end{description}

\end{fulllineitems}


\end{fulllineitems}

\index{printErrorMessage() (in module xmmPipeline.epicObj)@\spxentry{printErrorMessage()}\spxextra{in module xmmPipeline.epicObj}}

\begin{fulllineitems}
\phantomsection\label{\detokenize{xmmPipeline:xmmPipeline.epicObj.printErrorMessage}}
\pysigstartsignatures
\pysiglinewithargsret{\sphinxcode{\sphinxupquote{xmmPipeline.epicObj.}}\sphinxbfcode{\sphinxupquote{printErrorMessage}}}{\emph{\DUrole{n}{message}}}{}
\pysigstopsignatures
\end{fulllineitems}

\index{strToBool() (in module xmmPipeline.epicObj)@\spxentry{strToBool()}\spxextra{in module xmmPipeline.epicObj}}

\begin{fulllineitems}
\phantomsection\label{\detokenize{xmmPipeline:xmmPipeline.epicObj.strToBool}}
\pysigstartsignatures
\pysiglinewithargsret{\sphinxcode{\sphinxupquote{xmmPipeline.epicObj.}}\sphinxbfcode{\sphinxupquote{strToBool}}}{\emph{\DUrole{n}{s}}, \emph{\DUrole{n}{inverse}\DUrole{o}{=}\DUrole{default_value}{False}}}{}
\pysigstopsignatures
\end{fulllineitems}



\subsection{xmmPipeline.epicPileup module}
\label{\detokenize{xmmPipeline:module-xmmPipeline.epicPileup}}\label{\detokenize{xmmPipeline:xmmpipeline-epicpileup-module}}\index{module@\spxentry{module}!xmmPipeline.epicPileup@\spxentry{xmmPipeline.epicPileup}}\index{xmmPipeline.epicPileup@\spxentry{xmmPipeline.epicPileup}!module@\spxentry{module}}
\sphinxAtStartPar
08 May 2021:
—
For clearing the EPIC Pile\sphinxhyphen{}up in early time obsIDs.
See: \sphinxtitleref{https://www.cosmos.esa.int/web/xmm\sphinxhyphen{}newton/sas\sphinxhyphen{}thread\sphinxhyphen{}epatplot}

\sphinxAtStartPar
This would be used to correct for the Pile\sphinxhyphen{}up seperately for individual
obsIDs. Thus, the \sphinxtitleref{workdir} is changed to {\color{red}\bfseries{}\textasciigrave{}}workdir\textasciigrave{}+’/work’ folder.

\sphinxAtStartPar
This procedure will give the Source Annulus radii for which Pile\sphinxhyphen{}up is
negligible for any particular obsIDs. Henceforth, it is this Source region
for which the Light Curve and the Spectra extraction should be done.
\begin{description}
\sphinxlineitem{So, I will have to:}\begin{description}
\sphinxlineitem{\sphinxhyphen{}\textgreater{} Modify \sphinxtitleref{xmmPipeline.py} so that it knows which obsIDs have the}
\sphinxAtStartPar
Pile\sphinxhyphen{}up issue and need the Source region to lie within an Annulus.

\sphinxlineitem{\sphinxhyphen{}\textgreater{} Modify \sphinxtitleref{spectra.py} so that the Source Spectra are extracted from}
\sphinxAtStartPar
within this Annulus.

\end{description}

\end{description}

\sphinxAtStartPar
See the notes of Google Doc for open questions.

\sphinxAtStartPar
09 May 2021:
—
\sphinxtitleref{epatplot} calculates two diagnostic numbers which may be used to assess
the presence of pile\sphinxhyphen{}up: In the absence of pile\sphinxhyphen{}up, the 0.5 \sphinxhyphen{} 2.0 keV (default range)
observed\sphinxhyphen{}to\sphinxhyphen{}model singles and doubles pattern fractions ratios should both be
consistent with 1.0 within statistical errors (1 sigma errors are given).
If pile\sphinxhyphen{}up is present, the singles ratio will be smaller than 1.0 and the
doubles ratio will be larger than 1.0.

\sphinxAtStartPar
13 May 2021:
—
Adding \sphinxtitleref{correctPileUp} function.
Some late\sphinxhyphen{}time obsIDs 0770981001, 0810200501 and 0810200701 also appear to be Piled\sphinxhyphen{}up!
Gotta get this discrepancy cleared up.
\index{epicPileup (class in xmmPipeline.epicPileup)@\spxentry{epicPileup}\spxextra{class in xmmPipeline.epicPileup}}

\begin{fulllineitems}
\phantomsection\label{\detokenize{xmmPipeline:xmmPipeline.epicPileup.epicPileup}}
\pysigstartsignatures
\pysiglinewithargsret{\sphinxbfcode{\sphinxupquote{class\DUrole{w}{  }}}\sphinxcode{\sphinxupquote{xmmPipeline.epicPileup.}}\sphinxbfcode{\sphinxupquote{epicPileup}}}{\emph{\DUrole{n}{ra}}, \emph{\DUrole{n}{dec}}, \emph{\DUrole{n}{workdir}}, \emph{\DUrole{n}{sas\_dir}}, \emph{\DUrole{n}{headas}}, \emph{\DUrole{n}{sas\_ccfpath}}, \emph{\DUrole{n}{srcCircRadius}\DUrole{o}{=}\DUrole{default_value}{None}}, \emph{\DUrole{n}{bkgCircRadius}\DUrole{o}{=}\DUrole{default_value}{None}}, \emph{\DUrole{n}{dSrcThreshold}\DUrole{o}{=}\DUrole{default_value}{70}}, \emph{\DUrole{n}{bkgCircGap}\DUrole{o}{=}\DUrole{default_value}{2.5}}, \emph{\DUrole{n}{edetectmode}\DUrole{o}{=}\DUrole{default_value}{\textquotesingle{}chain\textquotesingle{}}}, \emph{\DUrole{n}{esp\_nsplinenodes}\DUrole{o}{=}\DUrole{default_value}{14}}, \emph{\DUrole{n}{gti\_indiThreshold}\DUrole{o}{=}\DUrole{default_value}{500}}, \emph{\DUrole{n}{gti\_combThreshold}\DUrole{o}{=}\DUrole{default_value}{None}}, \emph{\DUrole{n}{lcBinSize}\DUrole{o}{=}\DUrole{default_value}{25}}, \emph{\DUrole{n}{binBkglc}\DUrole{o}{=}\DUrole{default_value}{\textquotesingle{}no\textquotesingle{}}}, \emph{\DUrole{n}{saveFig}\DUrole{o}{=}\DUrole{default_value}{True}}, \emph{\DUrole{n}{showFig}\DUrole{o}{=}\DUrole{default_value}{True}}, \emph{\DUrole{n}{ignorePileup}\DUrole{o}{=}\DUrole{default_value}{False}}}{}
\pysigstopsignatures
\sphinxAtStartPar
Bases: \sphinxcode{\sphinxupquote{epicObj}}
\index{checkPileUp() (xmmPipeline.epicPileup.epicPileup method)@\spxentry{checkPileUp()}\spxextra{xmmPipeline.epicPileup.epicPileup method}}

\begin{fulllineitems}
\phantomsection\label{\detokenize{xmmPipeline:xmmPipeline.epicPileup.epicPileup.checkPileUp}}
\pysigstartsignatures
\pysiglinewithargsret{\sphinxbfcode{\sphinxupquote{checkPileUp}}}{}{}
\pysigstopsignatures
\sphinxAtStartPar
Uses \sphinxtitleref{epatplot} to check for the Pile\sphinxhyphen{}up issue.

\sphinxAtStartPar
First a GTI filtered Source Event list is extracted out of \sphinxtitleref{PN\_CCD\#.evts}
and then \sphinxtitleref{epatplot} is run to check the observed Spectra with the expected
pattern distribution function.

\sphinxAtStartPar
\sphinxtitleref{epatplot} calculates two diagnostic numbers which may be used to assess
the presence of pile\sphinxhyphen{}up: In the absence of pile\sphinxhyphen{}up, the 0.5 \sphinxhyphen{} 2.0 keV (default)
observed\sphinxhyphen{}to\sphinxhyphen{}model singles and doubles pattern fractions ratios should both be
consistent with 1.0 within statistical errors (1 sigma errors are given).
If pile\sphinxhyphen{}up is present, the singles ratio will be smaller than 1.0 and the
doubles ratio will be larger than 1.0.
See: \sphinxtitleref{https://www.cosmos.esa.int/web/xmm\sphinxhyphen{}newton/sas\sphinxhyphen{}thread\sphinxhyphen{}epatplot}

\sphinxAtStartPar
Note that \sphinxtitleref{pnGTI.fits} GTI file is used instead of \sphinxtitleref{combinedGTI\_obsID.fits}
because the latter results in very few useful data points.

\sphinxAtStartPar
If the obsID is found to be Piled\sphinxhyphen{}up then the user should run the \sphinxtitleref{correctPileUp}
function to correct for it.
\begin{quote}\begin{description}
\sphinxlineitem{Input}
\sphinxAtStartPar
Concatenated and Calibrated Event list, pnGTI file,
\sphinxtitleref{sourceCCDs} and the Source location parameters.

\sphinxlineitem{Output}
\sphinxAtStartPar
The Source Annulus radii together with Filtered Event List
\sphinxtitleref{source\_PN\_filtered\_obsID.evts} and the corresponding image
\sphinxtitleref{source\_PN\_filteredPattern.ps}.

\sphinxlineitem{Note}
\sphinxAtStartPar
The term \sphinxtitleref{filtered} means that the GTI have been applied.

\end{description}\end{quote}

\end{fulllineitems}

\index{correctPileUp() (xmmPipeline.epicPileup.epicPileup method)@\spxentry{correctPileUp()}\spxextra{xmmPipeline.epicPileup.epicPileup method}}

\begin{fulllineitems}
\phantomsection\label{\detokenize{xmmPipeline:xmmPipeline.epicPileup.epicPileup.correctPileUp}}
\pysigstartsignatures
\pysiglinewithargsret{\sphinxbfcode{\sphinxupquote{correctPileUp}}}{\emph{\DUrole{n}{srcRin}\DUrole{o}{=}\DUrole{default_value}{None}}}{}
\pysigstopsignatures
\sphinxAtStartPar
Uses \sphinxtitleref{epatplot} to correct and check for the Pile\sphinxhyphen{}up issue by removing
some pixels from the center of the Source circle.

\sphinxAtStartPar
The observation is corrected for Pile\sphinxhyphen{}up by removing some pixels from
the center of the Source. Now, this checking for Pile\sphinxhyphen{}up need be done
only once, for PN images, since, it will be the same for MOS, probably.

\sphinxAtStartPar
So, if Pile\sphinxhyphen{}up is found, one has to manually check for by iteratively
removing more and more center pixels and find when the Pile\sphinxhyphen{}up becomes
negligible. Then the Source circle would be an Annulus for this region.
The same Annulus parameters can be repeated for MOS1\&2.
See: \sphinxtitleref{https://www.cosmos.esa.int/web/xmm\sphinxhyphen{}newton/sas\sphinxhyphen{}thread\sphinxhyphen{}epatplot}

\sphinxAtStartPar
Note that \sphinxtitleref{pnGTI.fits} GTI file is used instead of \sphinxtitleref{combinedGTI\_obsID.fits}
because the latter results in very few useful data points.

\sphinxAtStartPar
The result of this procedure would be the radii values of the Source
Annulus for which Pile\sphinxhyphen{}up is negligible. These new found \sphinxtitleref{rIn} and \sphinxtitleref{rOut}
are then saved into to the \sphinxtitleref{epicObj} and are written to \sphinxtitleref{ccd\_coords\_info.pickle}
file. Henceforth, for these Pile\sphinxhyphen{}up obsIDs the Source region defined
from \sphinxtitleref{rIn} to \sphinxtitleref{rOut} will be used for obtaing the Spectra.
\begin{quote}\begin{description}
\sphinxlineitem{Input}
\sphinxAtStartPar
Concatenated and Calibrated Event list, pnGTI file,
\sphinxtitleref{sourceCCDs} and the Source location parameters.

\sphinxlineitem{Output}
\sphinxAtStartPar
The Source Annulus radii together with Filtered Event List
\sphinxtitleref{source\_PN\_filtered\_obsID.evts} and the corresponding image
\sphinxtitleref{source\_PN\_filteredPattern.ps}, and the Annulus files
\sphinxtitleref{source\_PN\_filteredAnnulus\_obsID.evts} and \sphinxtitleref{source\_PN\_filteredPattern\_Annulus.ps}
showing resolution of the Pile\sphinxhyphen{}up issue.

\sphinxlineitem{Note}
\sphinxAtStartPar
The term \sphinxtitleref{filtered} means that the GTI have been applied.

\end{description}\end{quote}

\end{fulllineitems}

\index{save\_pileUpResults() (xmmPipeline.epicPileup.epicPileup method)@\spxentry{save\_pileUpResults()}\spxextra{xmmPipeline.epicPileup.epicPileup method}}

\begin{fulllineitems}
\phantomsection\label{\detokenize{xmmPipeline:xmmPipeline.epicPileup.epicPileup.save_pileUpResults}}
\pysigstartsignatures
\pysiglinewithargsret{\sphinxbfcode{\sphinxupquote{save\_pileUpResults}}}{}{}
\pysigstopsignatures
\sphinxAtStartPar
Saves the \sphinxtitleref{*filteredPattern.ps} and \sphinxtitleref{*filteredAnnulus*.ps}
to the results directory.

\end{fulllineitems}


\end{fulllineitems}

\index{main() (in module xmmPipeline.epicPileup)@\spxentry{main()}\spxextra{in module xmmPipeline.epicPileup}}

\begin{fulllineitems}
\phantomsection\label{\detokenize{xmmPipeline:xmmPipeline.epicPileup.main}}
\pysigstartsignatures
\pysiglinewithargsret{\sphinxcode{\sphinxupquote{xmmPipeline.epicPileup.}}\sphinxbfcode{\sphinxupquote{main}}}{\emph{\DUrole{n}{args}}}{}
\pysigstopsignatures
\end{fulllineitems}



\subsection{xmmPipeline.epicPipeline module}
\label{\detokenize{xmmPipeline:module-xmmPipeline.epicPipeline}}\label{\detokenize{xmmPipeline:xmmpipeline-epicpipeline-module}}\index{module@\spxentry{module}!xmmPipeline.epicPipeline@\spxentry{xmmPipeline.epicPipeline}}\index{xmmPipeline.epicPipeline@\spxentry{xmmPipeline.epicPipeline}!module@\spxentry{module}}
\sphinxAtStartPar
10th May 2021:
—
Am making several changes to the code procedure arrangements.
See the GitHub repo, the Notes on Google Docs and the documentation for more information.
For previous docstring comments, see earlier code files, namely \sphinxtitleref{xmmPipeline.py}
\index{combineAndFind\_method() (in module xmmPipeline.epicPipeline)@\spxentry{combineAndFind\_method()}\spxextra{in module xmmPipeline.epicPipeline}}

\begin{fulllineitems}
\phantomsection\label{\detokenize{xmmPipeline:xmmPipeline.epicPipeline.combineAndFind_method}}
\pysigstartsignatures
\pysiglinewithargsret{\sphinxcode{\sphinxupquote{xmmPipeline.epicPipeline.}}\sphinxbfcode{\sphinxupquote{combineAndFind\_method}}}{\emph{\DUrole{n}{args}}}{}
\pysigstopsignatures
\end{fulllineitems}

\index{extractProds\_method() (in module xmmPipeline.epicPipeline)@\spxentry{extractProds\_method()}\spxextra{in module xmmPipeline.epicPipeline}}

\begin{fulllineitems}
\phantomsection\label{\detokenize{xmmPipeline:xmmPipeline.epicPipeline.extractProds_method}}
\pysigstartsignatures
\pysiglinewithargsret{\sphinxcode{\sphinxupquote{xmmPipeline.epicPipeline.}}\sphinxbfcode{\sphinxupquote{extractProds\_method}}}{\emph{\DUrole{n}{args}}}{}
\pysigstopsignatures
\end{fulllineitems}

\index{reduceData\_method() (in module xmmPipeline.epicPipeline)@\spxentry{reduceData\_method()}\spxextra{in module xmmPipeline.epicPipeline}}

\begin{fulllineitems}
\phantomsection\label{\detokenize{xmmPipeline:xmmPipeline.epicPipeline.reduceData_method}}
\pysigstartsignatures
\pysiglinewithargsret{\sphinxcode{\sphinxupquote{xmmPipeline.epicPipeline.}}\sphinxbfcode{\sphinxupquote{reduceData\_method}}}{\emph{\DUrole{n}{args}}}{}
\pysigstopsignatures
\end{fulllineitems}

\index{runIndiBkgCircFuncs\_method() (in module xmmPipeline.epicPipeline)@\spxentry{runIndiBkgCircFuncs\_method()}\spxextra{in module xmmPipeline.epicPipeline}}

\begin{fulllineitems}
\phantomsection\label{\detokenize{xmmPipeline:xmmPipeline.epicPipeline.runIndiBkgCircFuncs_method}}
\pysigstartsignatures
\pysiglinewithargsret{\sphinxcode{\sphinxupquote{xmmPipeline.epicPipeline.}}\sphinxbfcode{\sphinxupquote{runIndiBkgCircFuncs\_method}}}{\emph{\DUrole{n}{args}}}{}
\pysigstopsignatures
\end{fulllineitems}



\subsection{xmmPipeline.epicSpectra module}
\label{\detokenize{xmmPipeline:module-xmmPipeline.epicSpectra}}\label{\detokenize{xmmPipeline:xmmpipeline-epicspectra-module}}\index{module@\spxentry{module}!xmmPipeline.epicSpectra@\spxentry{xmmPipeline.epicSpectra}}\index{xmmPipeline.epicSpectra@\spxentry{xmmPipeline.epicSpectra}!module@\spxentry{module}}
\sphinxAtStartPar
10th May 2021:
—
Am making several changes to the code procedure arrangements.
See the GitHub repo, the Notes on Google Docs and the documentation for more information.
For previous docstring comments, see earlier code files, namely \sphinxtitleref{spectra.py}

\sphinxAtStartPar
11th May 2021:
—
Since model fitting for individual obsIDs has to be done manually, I don’t think
\sphinxtitleref{xspec\_fitSpectra} would be used now. \sphinxtitleref{pyXspec.py} will have to be run manually
for each obsIDs with the different model parameters varrying.

\sphinxAtStartPar
20th May 2021:
—
Completing the pending work of fitting Models to the Spectra.
Two criterions to be checked for each obsIDs:
\begin{itemize}
\item {} 
\sphinxAtStartPar
whether obsID is Piled\sphinxhyphen{}up? DONE!

\item {} 
\sphinxAtStartPar
whether obsID is in Small\sphinxhyphen{}mode? DONE!

\end{itemize}

\sphinxAtStartPar
21st May 2021:
—
Changed to using Flare Background filtered Event Lists for Spectra extraction.
Found a way to save the output parameter values from \sphinxtitleref{pyXspec.py}.
Completing the SpecModel fitting now.

\sphinxAtStartPar
13th June 2021:
—
Completing the work for obtaining individual instrument bkgCircs.
\begin{itemize}
\item {} 
\sphinxAtStartPar
Adding \sphinxtitleref{runIndiBkgCircFuncs} to the main function.

\end{itemize}

\sphinxAtStartPar
17th June 2021:
—
Adding features to read individual instrument bkgCircs from ccdCoordsPickle file.
And finally starting the manual specModel fits.
\index{epicSpectra (class in xmmPipeline.epicSpectra)@\spxentry{epicSpectra}\spxextra{class in xmmPipeline.epicSpectra}}

\begin{fulllineitems}
\phantomsection\label{\detokenize{xmmPipeline:xmmPipeline.epicSpectra.epicSpectra}}
\pysigstartsignatures
\pysiglinewithargsret{\sphinxbfcode{\sphinxupquote{class\DUrole{w}{  }}}\sphinxcode{\sphinxupquote{xmmPipeline.epicSpectra.}}\sphinxbfcode{\sphinxupquote{epicSpectra}}}{\emph{\DUrole{n}{doNotOverwriteModel}\DUrole{o}{=}\DUrole{default_value}{False}}, \emph{\DUrole{o}{*}\DUrole{n}{args}}, \emph{\DUrole{o}{**}\DUrole{n}{kwargs}}}{}
\pysigstopsignatures
\sphinxAtStartPar
Bases: \sphinxcode{\sphinxupquote{epicObj}}
\index{extractSpectra\_imageMode() (xmmPipeline.epicSpectra.epicSpectra method)@\spxentry{extractSpectra\_imageMode()}\spxextra{xmmPipeline.epicSpectra.epicSpectra method}}

\begin{fulllineitems}
\phantomsection\label{\detokenize{xmmPipeline:xmmPipeline.epicSpectra.epicSpectra.extractSpectra_imageMode}}
\pysigstartsignatures
\pysiglinewithargsret{\sphinxbfcode{\sphinxupquote{extractSpectra\_imageMode}}}{}{}
\pysigstopsignatures
\sphinxAtStartPar
Extracts the MOS12 and PN spectra in Image Mode.
See: \sphinxurl{https://www.cosmos.esa.int/web/xmm-newton/sas-thread-mos-spectrum},
\begin{quote}

\sphinxAtStartPar
\sphinxurl{https://www.cosmos.esa.int/web/xmm-newton/sas-thread-pn-spectrum}
\end{quote}
\begin{quote}\begin{description}
\sphinxlineitem{Input}
\sphinxAtStartPar
Flared Background filtered PN, MOS1\&2 Event Lists,
\sphinxtitleref{PNclean.ds}, \sphinxtitleref{MOS1clean.ds} and \sphinxtitleref{MOS2clean.ds}.
Alongwith SrcBkg coordinates and SourceCCDs resulting from \sphinxtitleref{epicPipeline.py}

\sphinxlineitem{Output}
\sphinxAtStartPar
MOS12 and PN SrcBkg Spectra in Image Mode, a Redistribution Matrix (\sphinxtitleref{rmf}) file and
an Effective Area Vector (\sphinxtitleref{arf}) file.

\end{description}\end{quote}
\subsubsection*{Notes}

\sphinxAtStartPar
Note that for Small Mode obsIDs, since the background is taken from PN image data,
instead of the overlap region data b’cuz MOS12 region was of small size, the background
region will lie outside of the image. Thus, only the Source Spectra can possibly be
extracted for Small Mode obsIDs. However, the ‘specgroup’ command requires both the
Source and the Background Spectra in order to create the grouped Spectrum. Therefore,
as of now, I am completely skipping getting the MOS12 Spectra for Small mode obsIDs.

\sphinxAtStartPar
20th May 2021: Now works with Pile\sphinxhyphen{}up corrected obsIDs.
Separate Source Spectrum file is generated for Piled\sphinxhyphen{}up cases,
\sphinxtitleref{inst\_spectrum\_source\_annulus\_obsID.fits}
All the other file names, including the Grouped Spectra is not altered.

\sphinxAtStartPar
17th June 2021: Adding feature for individual instrument bkgCircs.

\sphinxAtStartPar
If the total area of indi inst bkgCircs is greater than that of the overlap bkgCircs,
or if the former is not available for any given obsID, the overlap bkgCircs
are used.

\end{fulllineitems}

\index{readUpdatedCCDcoordsPickle() (xmmPipeline.epicSpectra.epicSpectra method)@\spxentry{readUpdatedCCDcoordsPickle()}\spxextra{xmmPipeline.epicSpectra.epicSpectra method}}

\begin{fulllineitems}
\phantomsection\label{\detokenize{xmmPipeline:xmmPipeline.epicSpectra.epicSpectra.readUpdatedCCDcoordsPickle}}
\pysigstartsignatures
\pysiglinewithargsret{\sphinxbfcode{\sphinxupquote{readUpdatedCCDcoordsPickle}}}{}{}
\pysigstopsignatures
\sphinxAtStartPar
Reads the individual instrument bkgCircs keywords from the updated
\sphinxtitleref{ccd\_coords\_info.pickle} file which contains the location parameters.

\end{fulllineitems}

\index{save\_spectraResults() (xmmPipeline.epicSpectra.epicSpectra method)@\spxentry{save\_spectraResults()}\spxextra{xmmPipeline.epicSpectra.epicSpectra method}}

\begin{fulllineitems}
\phantomsection\label{\detokenize{xmmPipeline:xmmPipeline.epicSpectra.epicSpectra.save_spectraResults}}
\pysigstartsignatures
\pysiglinewithargsret{\sphinxbfcode{\sphinxupquote{save\_spectraResults}}}{}{}
\pysigstopsignatures\begin{description}
\sphinxlineitem{Does one tasks}
\sphinxAtStartPar
Copies the Spectra images from each obsID directory to a results
folder in the main directory.

\end{description}

\end{fulllineitems}

\index{xspec\_fitSpectra() (xmmPipeline.epicSpectra.epicSpectra method)@\spxentry{xspec\_fitSpectra()}\spxextra{xmmPipeline.epicSpectra.epicSpectra method}}

\begin{fulllineitems}
\phantomsection\label{\detokenize{xmmPipeline:xmmPipeline.epicSpectra.epicSpectra.xspec_fitSpectra}}
\pysigstartsignatures
\pysiglinewithargsret{\sphinxbfcode{\sphinxupquote{xspec\_fitSpectra}}}{\emph{\DUrole{n}{instName}\DUrole{o}{=}\DUrole{default_value}{\textquotesingle{}PN\textquotesingle{}}}}{}
\pysigstopsignatures
\sphinxAtStartPar
Fits a model Spectra to the extracted Spectra data using \sphinxtitleref{pyXspec} package.
Since, \sphinxtitleref{pyXspec} requires HEA initialisation, a separate routine \sphinxtitleref{pyXspec.py}
is required to run it.
See: \sphinxurl{https://www.cosmos.esa.int/web/xmm-newton/sas-thread-xspec}
\begin{quote}\begin{description}
\sphinxlineitem{Input}
\sphinxAtStartPar
Extracted Spectra FITS file and input model parameters.

\sphinxlineitem{Output}
\sphinxAtStartPar
Spectra plots and \sphinxtitleref{specModelParams.json} file.

\end{description}\end{quote}

\sphinxAtStartPar
20th May 2021: Adding feature to allow model parameter value inputs.

\end{fulllineitems}


\end{fulllineitems}

\index{main() (in module xmmPipeline.epicSpectra)@\spxentry{main()}\spxextra{in module xmmPipeline.epicSpectra}}

\begin{fulllineitems}
\phantomsection\label{\detokenize{xmmPipeline:xmmPipeline.epicSpectra.main}}
\pysigstartsignatures
\pysiglinewithargsret{\sphinxcode{\sphinxupquote{xmmPipeline.epicSpectra.}}\sphinxbfcode{\sphinxupquote{main}}}{\emph{\DUrole{n}{args}}}{}
\pysigstopsignatures
\end{fulllineitems}



\subsection{xmmPipeline.findOverlap module}
\label{\detokenize{xmmPipeline:module-xmmPipeline.findOverlap}}\label{\detokenize{xmmPipeline:xmmpipeline-findoverlap-module}}\index{module@\spxentry{module}!xmmPipeline.findOverlap@\spxentry{xmmPipeline.findOverlap}}\index{xmmPipeline.findOverlap@\spxentry{xmmPipeline.findOverlap}!module@\spxentry{module}}
\sphinxAtStartPar
08\sphinxhyphen{}12 February 2021:
—
New Python function to obtain the overlap region.

\sphinxAtStartPar
22 February 2021:
—
Completing the modifications to the codes that were discussed on 13th of this month.

\sphinxAtStartPar
26 February 2021:
—
Had the call with \sphinxhref{mailto:Dheeraj@MIT}{Dheeraj@MIT} yester. Finalizing the things about this pipeline now.

\sphinxAtStartPar
10 March 2021:
—
The radius of the Source and Background circles given as input to evselect have to be in Sky Coords
or Pixel units and not arcsec, as I had been giving. Thus, the returned values from findOverlap.py
have to be altered.

\sphinxAtStartPar
30 March 2021:
—
Was thinking of using numba to accelerate the code, but since most of the functions
use matplotlib it is perhaps unfeasible.

\sphinxAtStartPar
29th April 2021:
—
The Background circles were too small in some cases.
Modifying the criterion for background circles radius and threshold.

\sphinxAtStartPar
Instead of taking some random point from with the points that are a threshold distance
away from other Sources, I can take the one with max radii. Currently, the Bkg radius is defined a\sphinxhyphen{}priori.

\sphinxAtStartPar
Tried using this, but there are so many complications.
So, have simply added another function to find the max bkg radii possible within the threshold,
to be used after the bkg coords have been found.

\sphinxAtStartPar
I think, I can have Bx1,By1 be found together with the xUseN,yUseN and continue having \sphinxtitleref{\_\_maxBkgRadius}
function for the Radii.
Yeah! This works. However, whether or not optimum radii are obtained, depends on which of the circles
is taken as first and which as the second.

\sphinxAtStartPar
I can also try to have the \sphinxtitleref{xmmObj} class as a meta\sphinxhyphen{}class for the \sphinxtitleref{findOverlap} class.
Nope! It’s better not to do this because I run \sphinxtitleref{findOverlap.py} for each obsID separately.

\sphinxAtStartPar
30th April 2021:
—
Made the second background point non\sphinxhyphen{}random as well.
The program takes a little longer now, but both the circles are large sized.
Gonna check if there’s any errors that come up because of this.
Also lowered the Bkg Circle Gap to 2.5 pixels instead of 5 pixles.

\sphinxAtStartPar
13th June 2021:
—
Writing the functions for obtaining individual background circles.
Learned about this awesome algorithm, PolyLabel for finding the \sphinxtitleref{pole of inaccessibility}
or the point of largest inscribed circle for any polygon. My problem of finding the Bkg
circles is exactly the same.
See: \sphinxurl{https://github.com/Twista/python-polylabel}

\sphinxAtStartPar
Damn! That polylabel requires the boundary of the polygon to be defined and hell,
I have the whole points in the region. Running the \sphinxcode{\sphinxupquote{polylabel(points)}} on that full
throttles the system. Hell!
My own solution is pretty elegant then. Better suited for my purpose.

\sphinxAtStartPar
Wait a minute, what I can do to find more number of Bkg points is first find the first
point and then remove all the points lying in the 2R distance from the center of this
first Bkg point and repeat. Let’s try this!

\sphinxAtStartPar
Yosh! This works brilliantly. Both the BkgPts are now randomly generated in the given
area satisfying the coditions.
Now, I can theoretically iterate a few times until I get both the BkgCircs of sufficiently
useful area. However, what my aim has been ever since the start was to have the algorithm
automatically find two of the largest BkgCircs satisfying the coditions.
Hhmm… I think that would be a pretty daunting task and a difficult one at that,
since I haven’t go just what I have been wanting since a long time now.

\sphinxAtStartPar
What is much better in my opinion is that I can repeat the procedure of filtering out the
last found BkgPt + 2R dist points n number of times to obtain n number of Bkg points.
Let’s try that out now!

\sphinxAtStartPar
This recursive way of finding the background circles takes care of overlap between the
circles all by itself, since the \sphinxcode{\sphinxupquote{xToUse}} and \sphinxcode{\sphinxupquote{yToUse}} are recursively filtered for
the circle area.
Yup! Have verified that most of the times, this recursive way works! Woohooo yayyy!
Such a great thing.
For the times when it doesn’t work, one of the bkgCirc overlaps with another bkgCircs,
and consquensly the radius value for comes out to be negative.
I am adding feature to raise a warning when this happens. A simple rerun of the code
will remove the problem and yield correct bkgCircs, although yeah! I may need to check
this up with all the obsIDs.
And it would be better to have the user input the \sphinxcode{\sphinxupquote{nBkgCirc}} value, so that if it’s
not possible to obtain more than 2 bkgCircs, the user can find 2 bkgCircs atleast.
\index{findOverlap (class in xmmPipeline.findOverlap)@\spxentry{findOverlap}\spxextra{class in xmmPipeline.findOverlap}}

\begin{fulllineitems}
\phantomsection\label{\detokenize{xmmPipeline:xmmPipeline.findOverlap.findOverlap}}
\pysigstartsignatures
\pysiglinewithargsret{\sphinxbfcode{\sphinxupquote{class\DUrole{w}{  }}}\sphinxcode{\sphinxupquote{xmmPipeline.findOverlap.}}\sphinxbfcode{\sphinxupquote{findOverlap}}}{\emph{\DUrole{n}{workdir}}, \emph{\DUrole{n}{obsID}\DUrole{o}{=}\DUrole{default_value}{\textquotesingle{}NA\textquotesingle{}}}, \emph{\DUrole{n}{pnCCD}\DUrole{o}{=}\DUrole{default_value}{\textquotesingle{}4\textquotesingle{}}}, \emph{\DUrole{n}{srcCoords}\DUrole{o}{=}\DUrole{default_value}{None}}, \emph{\DUrole{n}{otherSrc}\DUrole{o}{=}\DUrole{default_value}{None}}, \emph{\DUrole{n}{gap}\DUrole{o}{=}\DUrole{default_value}{2.5}}, \emph{\DUrole{n}{srcR}\DUrole{o}{=}\DUrole{default_value}{None}}, \emph{\DUrole{n}{bkgR}\DUrole{o}{=}\DUrole{default_value}{None}}, \emph{\DUrole{n}{srcThreshold}\DUrole{o}{=}\DUrole{default_value}{None}}, \emph{\DUrole{n}{saveFig}\DUrole{o}{=}\DUrole{default_value}{True}}, \emph{\DUrole{n}{showFig}\DUrole{o}{=}\DUrole{default_value}{True}}}{}
\pysigstopsignatures
\sphinxAtStartPar
Bases: \sphinxhref{https://docs.python.org/3/library/functions.html\#object}{\sphinxcode{\sphinxupquote{object}}}
\index{backgroundCircles() (xmmPipeline.findOverlap.findOverlap method)@\spxentry{backgroundCircles()}\spxextra{xmmPipeline.findOverlap.findOverlap method}}

\begin{fulllineitems}
\phantomsection\label{\detokenize{xmmPipeline:xmmPipeline.findOverlap.findOverlap.backgroundCircles}}
\pysigstartsignatures
\pysiglinewithargsret{\sphinxbfcode{\sphinxupquote{backgroundCircles}}}{\emph{\DUrole{n}{axes}}}{}
\pysigstopsignatures\begin{quote}\begin{description}
\sphinxlineitem{Input}
\sphinxAtStartPar
The output from ‘obtainOverlap()’ function aka detected overlap region.,
and the Source coordinates.

\sphinxlineitem{Output}
\sphinxAtStartPar
Coordinates and Radii of the Background Circles.

\end{description}\end{quote}

\end{fulllineitems}

\index{backgroundCircles\_indi() (xmmPipeline.findOverlap.findOverlap method)@\spxentry{backgroundCircles\_indi()}\spxextra{xmmPipeline.findOverlap.findOverlap method}}

\begin{fulllineitems}
\phantomsection\label{\detokenize{xmmPipeline:xmmPipeline.findOverlap.findOverlap.backgroundCircles_indi}}
\pysigstartsignatures
\pysiglinewithargsret{\sphinxbfcode{\sphinxupquote{backgroundCircles\_indi}}}{\emph{\DUrole{n}{ax}}, \emph{\DUrole{n}{inst}}, \emph{\DUrole{n}{nBkgCircs}\DUrole{o}{=}\DUrole{default_value}{3}}}{}
\pysigstopsignatures
\sphinxAtStartPar
Automatically finds Background Circles in the instrument data passed.
Modified the function a great deal from the earlier naive \sphinxtitleref{backgroundCircles} function.
\begin{quote}\begin{description}
\sphinxlineitem{Input}
\sphinxAtStartPar
\sphinxtitleref{inst\_CCD\#\#\_image.fits} files containing the extracted image data of the instrument.
\sphinxcode{\sphinxupquote{srcCoords}} dictionary containing coords of the main Source.
\sphinxcode{\sphinxupquote{otherSrc}} dictionary containing coords of other Sources.

\sphinxlineitem{Requires}
\sphinxAtStartPar
Call to \sphinxtitleref{epicObj.find\_otherSources\_indi} function.

\sphinxlineitem{Parameters}\begin{itemize}
\item {} 
\sphinxAtStartPar
\sphinxstyleliteralstrong{\sphinxupquote{ax}} (\sphinxstyleliteralemphasis{\sphinxupquote{obj}}) \textendash{} \sphinxtitleref{matplotlib.pyplot.subplots} object for plotting the images.

\item {} 
\sphinxAtStartPar
\sphinxstyleliteralstrong{\sphinxupquote{inst}} (\sphinxhref{https://docs.python.org/3/library/stdtypes.html\#str}{\sphinxstyleliteralemphasis{\sphinxupquote{str}}}) \textendash{} name of the instrument of use.

\item {} 
\sphinxAtStartPar
\sphinxstyleliteralstrong{\sphinxupquote{nBkgCircs}} (\sphinxhref{https://docs.python.org/3/library/functions.html\#int}{\sphinxstyleliteralemphasis{\sphinxupquote{int}}}) \textendash{} number of background circles to detect. Defaults to 3. (Optional)

\end{itemize}

\sphinxlineitem{Output}
\sphinxAtStartPar
Coordinates and Radii of the Background Circles.

\end{description}\end{quote}
\subsubsection*{Notes}

\sphinxAtStartPar
\sphinxcode{\sphinxupquote{small\sphinxhyphen{}mode}} MOS checking is invalid here since overlap is not been detected.

\end{fulllineitems}

\index{createOutputImages() (xmmPipeline.findOverlap.findOverlap method)@\spxentry{createOutputImages()}\spxextra{xmmPipeline.findOverlap.findOverlap method}}

\begin{fulllineitems}
\phantomsection\label{\detokenize{xmmPipeline:xmmPipeline.findOverlap.findOverlap.createOutputImages}}
\pysigstartsignatures
\pysiglinewithargsret{\sphinxbfcode{\sphinxupquote{createOutputImages}}}{\emph{\DUrole{n}{inst}\DUrole{o}{=}\DUrole{default_value}{None}}}{}
\pysigstopsignatures
\sphinxAtStartPar
Creates the \sphinxtitleref{source\_background\_circles.png} output image showing
the other Sources and the Source and Background circles overlayed
on top of the image.
The overlap (‘PNMOS12\_image\_full’ for Small\sphinxhyphen{}Mode cases) is used
as the backdrop for cases when the overlap detection is performed
and the cicles are found using that overlap data.
Otherwise, \sphinxtitleref{inst\_CCD\#\#\_image} is used as the backdrop.
\begin{quote}\begin{description}
\sphinxlineitem{Parameters}
\sphinxAtStartPar
\sphinxstyleliteralstrong{\sphinxupquote{inst}} (\sphinxhref{https://docs.python.org/3/library/stdtypes.html\#str}{\sphinxstyleliteralemphasis{\sphinxupquote{str}}}) \textendash{} Name of the instrument to use. Defaults to None.
If None then it is assumed at the overlap detected is
performed.

\sphinxlineitem{Input}
\sphinxAtStartPar
findOverlap object with all previous functions already run.

\sphinxlineitem{Output}
\sphinxAtStartPar
Output image saved in the workdir of the obsID.
Image file name is \sphinxtitleref{source\_background\_circles.png} for the
default case and \sphinxtitleref{srcBkg\_circs\_inst.png} for when individual
instrument background circles are been found.

\end{description}\end{quote}

\end{fulllineitems}

\index{main() (xmmPipeline.findOverlap.findOverlap method)@\spxentry{main()}\spxextra{xmmPipeline.findOverlap.findOverlap method}}

\begin{fulllineitems}
\phantomsection\label{\detokenize{xmmPipeline:xmmPipeline.findOverlap.findOverlap.main}}
\pysigstartsignatures
\pysiglinewithargsret{\sphinxbfcode{\sphinxupquote{main}}}{\emph{\DUrole{n}{inst}\DUrole{o}{=}\DUrole{default_value}{None}}}{}
\pysigstopsignatures
\sphinxAtStartPar
The main function to sequentially run the functions in \sphinxtitleref{findOverlap} class.
\begin{quote}\begin{description}
\sphinxlineitem{Returns}
\sphinxAtStartPar

\sphinxAtStartPar
(x, y, r) coords and radius of first background circle.
bCircle2 (tuple): (x, y, r) coords and radius of second background circle.
\begin{description}
\sphinxlineitem{correctSrc (tuple): corrected (x, y) coords of the main Source.}
\sphinxAtStartPar
The original main Source coords found by the \sphinxtitleref{epicObj.findSourceCCD}
function using the XMMSAS \sphinxcode{\sphinxupquote{ecoordconv}} utility are a little bit shifted
from what the main Source likely is as found from \sphinxcode{\sphinxupquote{edetect\_chain}} while
finding other Sources. So, the main Source coords are changed to the coords
of the nearest Source to it in the \sphinxcode{\sphinxupquote{otherSrc}} dictionary.

\sphinxlineitem{srcR (float): radius of the Source circle or outer radius of the Source Annulus}
\sphinxAtStartPar
for the Piled\sphinxhyphen{}up cases, although a full circle is detected here.

\sphinxlineitem{isSmallMode (bool): True, if absolute difference between flattened MOS12 and}
\sphinxAtStartPar
PNMOS12 image arrays is less than 500. False otherwise.
See issue \#28 in the xmmPipeline repository in this regard,
\sphinxurl{https://github.com/Suyog7130/xmmPipeline/issues/28}

\end{description}


\sphinxlineitem{Return type}
\sphinxAtStartPar
bCircle1 (\sphinxhref{https://docs.python.org/3/library/stdtypes.html\#tuple}{tuple})

\end{description}\end{quote}
\subsubsection*{Notes}

\sphinxAtStartPar
All returned coordinates are in \sphinxstylestrong{Physical Units}.
These (x, y) are termed as (\sphinxstylestrong{RAWX}, \sphinxstylestrong{RAWY}) in DS9.
The radius values are in \sphinxstylestrong{Sky Coordinates}, which I think is the name XMMSAS gives
to \sphinxstylestrong{Physical Units}. Each Sky Coord unit equals \sphinxstylestrong{0.05 arcsec}.
So, a radius of 640 equals 32 arcsec.
See issue \#24 for more details, \sphinxurl{https://github.com/Suyog7130/xmmPipeline/issues/24}

\end{fulllineitems}

\index{obtainOverlap() (xmmPipeline.findOverlap.findOverlap method)@\spxentry{obtainOverlap()}\spxextra{xmmPipeline.findOverlap.findOverlap method}}

\begin{fulllineitems}
\phantomsection\label{\detokenize{xmmPipeline:xmmPipeline.findOverlap.findOverlap.obtainOverlap}}
\pysigstartsignatures
\pysiglinewithargsret{\sphinxbfcode{\sphinxupquote{obtainOverlap}}}{\emph{\DUrole{n}{axes}}}{}
\pysigstopsignatures\begin{quote}\begin{description}
\sphinxlineitem{Input}
\sphinxAtStartPar
The workdir where the PN, MOS and PNMOS12 FITS files are located.

\sphinxlineitem{Output}
\sphinxAtStartPar
The detected overlap region if it is larger in size than then MOS region,
or else the PNMOS12 image data.

\end{description}\end{quote}

\end{fulllineitems}


\end{fulllineitems}

\index{imageMOS12() (in module xmmPipeline.findOverlap)@\spxentry{imageMOS12()}\spxextra{in module xmmPipeline.findOverlap}}

\begin{fulllineitems}
\phantomsection\label{\detokenize{xmmPipeline:xmmPipeline.findOverlap.imageMOS12}}
\pysigstartsignatures
\pysiglinewithargsret{\sphinxcode{\sphinxupquote{xmmPipeline.findOverlap.}}\sphinxbfcode{\sphinxupquote{imageMOS12}}}{\emph{\DUrole{n}{workdir}}}{}
\pysigstopsignatures
\end{fulllineitems}

\index{imageMOS12\_soft() (in module xmmPipeline.findOverlap)@\spxentry{imageMOS12\_soft()}\spxextra{in module xmmPipeline.findOverlap}}

\begin{fulllineitems}
\phantomsection\label{\detokenize{xmmPipeline:xmmPipeline.findOverlap.imageMOS12_soft}}
\pysigstartsignatures
\pysiglinewithargsret{\sphinxcode{\sphinxupquote{xmmPipeline.findOverlap.}}\sphinxbfcode{\sphinxupquote{imageMOS12\_soft}}}{\emph{\DUrole{n}{workdir}}}{}
\pysigstopsignatures
\end{fulllineitems}

\index{imagePN() (in module xmmPipeline.findOverlap)@\spxentry{imagePN()}\spxextra{in module xmmPipeline.findOverlap}}

\begin{fulllineitems}
\phantomsection\label{\detokenize{xmmPipeline:xmmPipeline.findOverlap.imagePN}}
\pysigstartsignatures
\pysiglinewithargsret{\sphinxcode{\sphinxupquote{xmmPipeline.findOverlap.}}\sphinxbfcode{\sphinxupquote{imagePN}}}{\emph{\DUrole{n}{workdir}}, \emph{\DUrole{n}{pnCCD}}}{}
\pysigstopsignatures
\end{fulllineitems}

\index{imagePNMOS12() (in module xmmPipeline.findOverlap)@\spxentry{imagePNMOS12()}\spxextra{in module xmmPipeline.findOverlap}}

\begin{fulllineitems}
\phantomsection\label{\detokenize{xmmPipeline:xmmPipeline.findOverlap.imagePNMOS12}}
\pysigstartsignatures
\pysiglinewithargsret{\sphinxcode{\sphinxupquote{xmmPipeline.findOverlap.}}\sphinxbfcode{\sphinxupquote{imagePNMOS12}}}{\emph{\DUrole{n}{workdir}}}{}
\pysigstopsignatures
\end{fulllineitems}

\index{imagePNMOS12\_soft() (in module xmmPipeline.findOverlap)@\spxentry{imagePNMOS12\_soft()}\spxextra{in module xmmPipeline.findOverlap}}

\begin{fulllineitems}
\phantomsection\label{\detokenize{xmmPipeline:xmmPipeline.findOverlap.imagePNMOS12_soft}}
\pysigstartsignatures
\pysiglinewithargsret{\sphinxcode{\sphinxupquote{xmmPipeline.findOverlap.}}\sphinxbfcode{\sphinxupquote{imagePNMOS12\_soft}}}{\emph{\DUrole{n}{workdir}}}{}
\pysigstopsignatures
\end{fulllineitems}

\index{imagePN\_soft() (in module xmmPipeline.findOverlap)@\spxentry{imagePN\_soft()}\spxextra{in module xmmPipeline.findOverlap}}

\begin{fulllineitems}
\phantomsection\label{\detokenize{xmmPipeline:xmmPipeline.findOverlap.imagePN_soft}}
\pysigstartsignatures
\pysiglinewithargsret{\sphinxcode{\sphinxupquote{xmmPipeline.findOverlap.}}\sphinxbfcode{\sphinxupquote{imagePN\_soft}}}{\emph{\DUrole{n}{workdir}}, \emph{\DUrole{n}{pnCCD}}}{}
\pysigstopsignatures
\end{fulllineitems}



\subsection{xmmPipeline.getSpectraFitsTable module}
\label{\detokenize{xmmPipeline:module-xmmPipeline.getSpectraFitsTable}}\label{\detokenize{xmmPipeline:xmmpipeline-getspectrafitstable-module}}\index{module@\spxentry{module}!xmmPipeline.getSpectraFitsTable@\spxentry{xmmPipeline.getSpectraFitsTable}}\index{xmmPipeline.getSpectraFitsTable@\spxentry{xmmPipeline.getSpectraFitsTable}!module@\spxentry{module}}
\sphinxAtStartPar
22nd May 2021:
—
Have completed much of the pending work regarding Spectral Analysis
of ASASSN\sphinxhyphen{}14li. This Python routine is to read the \sphinxtitleref{specModelParams.json}
files for each of the obsIDs and create a \sphinxtitleref{specResultsTable} using them.

\sphinxAtStartPar
I don’t the obsIDs thing will be valid here.
Why would one not wanna include some obsIDs in the final table?
Anyway.

\sphinxAtStartPar
23rd May 2021:
—
The primary \sphinxcode{\sphinxupquote{for}} loop should be for model since each model will have a
separate \sphinxcode{\sphinxupquote{specResultsTable}}.

\sphinxAtStartPar
18th June 2021:
—
Finishing up the manual specModel fitting.

\sphinxAtStartPar
NOTE: Be wary of the specModel name with single additive model component
when it has parentheses around it. The model name saved in the JSON file
and that been given as input, should exactly match.
See issue \#31 in regard to this. \sphinxurl{https://github.com/Suyog7130/xmmPipeline/issues/31}

\sphinxAtStartPar
Swaping the Row of the \sphinxtitleref{specResultsTable} with the columns will take sometime.
\index{main() (in module xmmPipeline.getSpectraFitsTable)@\spxentry{main()}\spxextra{in module xmmPipeline.getSpectraFitsTable}}

\begin{fulllineitems}
\phantomsection\label{\detokenize{xmmPipeline:xmmPipeline.getSpectraFitsTable.main}}
\pysigstartsignatures
\pysiglinewithargsret{\sphinxcode{\sphinxupquote{xmmPipeline.getSpectraFitsTable.}}\sphinxbfcode{\sphinxupquote{main}}}{\emph{\DUrole{n}{args}}}{}
\pysigstopsignatures
\end{fulllineitems}



\subsection{xmmPipeline.plotAnal module}
\label{\detokenize{xmmPipeline:module-xmmPipeline.plotAnal}}\label{\detokenize{xmmPipeline:xmmpipeline-plotanal-module}}\index{module@\spxentry{module}!xmmPipeline.plotAnal@\spxentry{xmmPipeline.plotAnal}}\index{xmmPipeline.plotAnal@\spxentry{xmmPipeline.plotAnal}!module@\spxentry{module}}
\sphinxAtStartPar
13th May 2020:
A module to ease up beatifying the plots.

\sphinxAtStartPar
13th June 2020:
Made changes on this.
If the fit is linear regression then the rSquared coefficient of determination
is the square of the pearson correlation coefficient.

\sphinxAtStartPar
20th August 2020:
Adding keyword for number of ticks shown.

\sphinxAtStartPar
23rd September 2020:
Adding keyword for axis colours via ax.spines{[}‘bottom’{]}.set\_color()
\begin{quote}

\sphinxAtStartPar
NOTE: spines can only be used with figure and not subplots!
\end{quote}

\sphinxAtStartPar
24th January 2021:
Made modifications in the beautifyPlot().

\sphinxAtStartPar
14th April 2021:
—
Adding modifications to the beautifyPlot function.
\index{plotAnal (class in xmmPipeline.plotAnal)@\spxentry{plotAnal}\spxextra{class in xmmPipeline.plotAnal}}

\begin{fulllineitems}
\phantomsection\label{\detokenize{xmmPipeline:xmmPipeline.plotAnal.plotAnal}}
\pysigstartsignatures
\pysiglinewithargsret{\sphinxbfcode{\sphinxupquote{class\DUrole{w}{  }}}\sphinxcode{\sphinxupquote{xmmPipeline.plotAnal.}}\sphinxbfcode{\sphinxupquote{plotAnal}}}{\emph{\DUrole{n}{figure}}, \emph{\DUrole{n}{x}}, \emph{\DUrole{n}{y}}}{}
\pysigstopsignatures
\sphinxAtStartPar
Bases: \sphinxhref{https://docs.python.org/3/library/functions.html\#object}{\sphinxcode{\sphinxupquote{object}}}
\index{PSNR() (xmmPipeline.plotAnal.plotAnal method)@\spxentry{PSNR()}\spxextra{xmmPipeline.plotAnal.plotAnal method}}

\begin{fulllineitems}
\phantomsection\label{\detokenize{xmmPipeline:xmmPipeline.plotAnal.plotAnal.PSNR}}
\pysigstartsignatures
\pysiglinewithargsret{\sphinxbfcode{\sphinxupquote{PSNR}}}{\emph{\DUrole{n}{x}}, \emph{\DUrole{n}{y}}, \emph{\DUrole{n}{xloc}\DUrole{o}{=}\DUrole{default_value}{0.05}}, \emph{\DUrole{n}{yloc}\DUrole{o}{=}\DUrole{default_value}{0.1}}, \emph{\DUrole{n}{decPlace}\DUrole{o}{=}\DUrole{default_value}{2}}, \emph{\DUrole{n}{fontsize}\DUrole{o}{=}\DUrole{default_value}{15}}}{}
\pysigstopsignatures
\sphinxAtStartPar
Calculates Peak Signal\sphinxhyphen{}to\sphinxhyphen{}Noise Ratio between x and y and plots its value on the figure.

\end{fulllineitems}

\index{SSIM() (xmmPipeline.plotAnal.plotAnal method)@\spxentry{SSIM()}\spxextra{xmmPipeline.plotAnal.plotAnal method}}

\begin{fulllineitems}
\phantomsection\label{\detokenize{xmmPipeline:xmmPipeline.plotAnal.plotAnal.SSIM}}
\pysigstartsignatures
\pysiglinewithargsret{\sphinxbfcode{\sphinxupquote{SSIM}}}{\emph{\DUrole{n}{x}}, \emph{\DUrole{n}{y}}, \emph{\DUrole{n}{xloc}\DUrole{o}{=}\DUrole{default_value}{0.05}}, \emph{\DUrole{n}{yloc}\DUrole{o}{=}\DUrole{default_value}{0.1}}, \emph{\DUrole{n}{decPlace}\DUrole{o}{=}\DUrole{default_value}{2}}, \emph{\DUrole{n}{fontsize}\DUrole{o}{=}\DUrole{default_value}{15}}, \emph{\DUrole{n}{scikit\_value}\DUrole{o}{=}\DUrole{default_value}{False}}}{}
\pysigstopsignatures
\sphinxAtStartPar
Calculates the Structural SIMilarity index between two image data given by x and y.
The calculated value is then plotted on the figure passed.

\end{fulllineitems}

\index{SSIM\_value() (xmmPipeline.plotAnal.plotAnal method)@\spxentry{SSIM\_value()}\spxextra{xmmPipeline.plotAnal.plotAnal method}}

\begin{fulllineitems}
\phantomsection\label{\detokenize{xmmPipeline:xmmPipeline.plotAnal.plotAnal.SSIM_value}}
\pysigstartsignatures
\pysiglinewithargsret{\sphinxbfcode{\sphinxupquote{SSIM\_value}}}{\emph{\DUrole{n}{y}}, \emph{\DUrole{n}{decPlace}\DUrole{o}{=}\DUrole{default_value}{2}}, \emph{\DUrole{n}{scikit\_value}\DUrole{o}{=}\DUrole{default_value}{False}}}{}
\pysigstopsignatures
\sphinxAtStartPar
Calculates the Structural SIMilarity index between two image data given by x and y.

\end{fulllineitems}

\index{beautifyPlot() (xmmPipeline.plotAnal.plotAnal method)@\spxentry{beautifyPlot()}\spxextra{xmmPipeline.plotAnal.plotAnal method}}

\begin{fulllineitems}
\phantomsection\label{\detokenize{xmmPipeline:xmmPipeline.plotAnal.plotAnal.beautifyPlot}}
\pysigstartsignatures
\pysiglinewithargsret{\sphinxbfcode{\sphinxupquote{beautifyPlot}}}{\emph{\DUrole{n}{labelsize}\DUrole{o}{=}\DUrole{default_value}{11}}, \emph{\DUrole{n}{lengthMajor}\DUrole{o}{=}\DUrole{default_value}{10}}, \emph{\DUrole{n}{tickNum}\DUrole{o}{=}\DUrole{default_value}{10}}, \emph{\DUrole{n}{tickDirection}\DUrole{o}{=}\DUrole{default_value}{\textquotesingle{}out\textquotesingle{}}}, \emph{\DUrole{n}{lengthMinor}\DUrole{o}{=}\DUrole{default_value}{5}}, \emph{\DUrole{n}{minor}\DUrole{o}{=}\DUrole{default_value}{False}}, \emph{\DUrole{n}{grid}\DUrole{o}{=}\DUrole{default_value}{False}}, \emph{\DUrole{n}{axisColor}\DUrole{o}{=}\DUrole{default_value}{None}}, \emph{\DUrole{n}{yTicks}\DUrole{o}{=}\DUrole{default_value}{True}}, \emph{\DUrole{n}{xTicks}\DUrole{o}{=}\DUrole{default_value}{True}}, \emph{\DUrole{n}{logFormat}\DUrole{o}{=}\DUrole{default_value}{None}}, \emph{\DUrole{n}{logXformat}\DUrole{o}{=}\DUrole{default_value}{None}}, \emph{\DUrole{n}{logYformat}\DUrole{o}{=}\DUrole{default_value}{None}}, \emph{\DUrole{n}{logMinor}\DUrole{o}{=}\DUrole{default_value}{True}}, \emph{\DUrole{n}{minorLabel}\DUrole{o}{=}\DUrole{default_value}{False}}, \emph{\DUrole{n}{logXminorLabel}\DUrole{o}{=}\DUrole{default_value}{False}}, \emph{\DUrole{n}{logYminorLabel}\DUrole{o}{=}\DUrole{default_value}{False}}}{}
\pysigstopsignatures
\sphinxAtStartPar
If axes subplots is used, with multiple rows and columns, then the ndarray of axes subobjects
can be flattened:
\begin{quote}

\sphinxAtStartPar
axes.flatten().tolist()
\end{quote}

\sphinxAtStartPar
The axisColor takes in a dictionary axisColor with keys as the spine position and values
as the colour that has to be inserted.
Be wary of the american spelling of ‘colour’, ‘color’, used in here.
NOTE: spines can only be used with figure and not subplots!

\end{fulllineitems}

\index{linearFit() (xmmPipeline.plotAnal.plotAnal method)@\spxentry{linearFit()}\spxextra{xmmPipeline.plotAnal.plotAnal method}}

\begin{fulllineitems}
\phantomsection\label{\detokenize{xmmPipeline:xmmPipeline.plotAnal.plotAnal.linearFit}}
\pysigstartsignatures
\pysiglinewithargsret{\sphinxbfcode{\sphinxupquote{linearFit}}}{\emph{\DUrole{n}{x}}, \emph{\DUrole{n}{y}}, \emph{\DUrole{n}{lineDetails}\DUrole{o}{=}\DUrole{default_value}{True}}, \emph{\DUrole{n}{xloc}\DUrole{o}{=}\DUrole{default_value}{0.05}}, \emph{\DUrole{n}{yloc}\DUrole{o}{=}\DUrole{default_value}{0.1}}, \emph{\DUrole{n}{fontsize}\DUrole{o}{=}\DUrole{default_value}{15}}, \emph{\DUrole{n}{color}\DUrole{o}{=}\DUrole{default_value}{\textquotesingle{}darkorange\textquotesingle{}}}}{}
\pysigstopsignatures
\end{fulllineitems}

\index{linearFit\_plotCal() (xmmPipeline.plotAnal.plotAnal method)@\spxentry{linearFit\_plotCal()}\spxextra{xmmPipeline.plotAnal.plotAnal method}}

\begin{fulllineitems}
\phantomsection\label{\detokenize{xmmPipeline:xmmPipeline.plotAnal.plotAnal.linearFit_plotCal}}
\pysigstartsignatures
\pysiglinewithargsret{\sphinxbfcode{\sphinxupquote{linearFit\_plotCal}}}{\emph{\DUrole{n}{x}}, \emph{\DUrole{n}{y}}, \emph{\DUrole{n}{xloc}\DUrole{o}{=}\DUrole{default_value}{0.05}}, \emph{\DUrole{n}{yloc}\DUrole{o}{=}\DUrole{default_value}{0.15}}, \emph{\DUrole{n}{fontsize}\DUrole{o}{=}\DUrole{default_value}{15}}}{}
\pysigstopsignatures
\end{fulllineitems}

\index{pearson() (xmmPipeline.plotAnal.plotAnal method)@\spxentry{pearson()}\spxextra{xmmPipeline.plotAnal.plotAnal method}}

\begin{fulllineitems}
\phantomsection\label{\detokenize{xmmPipeline:xmmPipeline.plotAnal.plotAnal.pearson}}
\pysigstartsignatures
\pysiglinewithargsret{\sphinxbfcode{\sphinxupquote{pearson}}}{\emph{\DUrole{n}{x}}, \emph{\DUrole{n}{y}}, \emph{\DUrole{n}{line}\DUrole{o}{=}\DUrole{default_value}{False}}, \emph{\DUrole{n}{xloc}\DUrole{o}{=}\DUrole{default_value}{0.05}}, \emph{\DUrole{n}{yloc}\DUrole{o}{=}\DUrole{default_value}{0.1}}, \emph{\DUrole{n}{decPlace}\DUrole{o}{=}\DUrole{default_value}{2}}, \emph{\DUrole{n}{rSquared}\DUrole{o}{=}\DUrole{default_value}{False}}, \emph{\DUrole{n}{fontsize}\DUrole{o}{=}\DUrole{default_value}{15}}, \emph{\DUrole{n}{linecolor}\DUrole{o}{=}\DUrole{default_value}{\textquotesingle{}darkorange\textquotesingle{}}}, \emph{\DUrole{n}{lineDetails}\DUrole{o}{=}\DUrole{default_value}{False}}}{}
\pysigstopsignatures
\end{fulllineitems}

\index{pearson\_spearman() (xmmPipeline.plotAnal.plotAnal method)@\spxentry{pearson\_spearman()}\spxextra{xmmPipeline.plotAnal.plotAnal method}}

\begin{fulllineitems}
\phantomsection\label{\detokenize{xmmPipeline:xmmPipeline.plotAnal.plotAnal.pearson_spearman}}
\pysigstartsignatures
\pysiglinewithargsret{\sphinxbfcode{\sphinxupquote{pearson\_spearman}}}{\emph{\DUrole{n}{x}}, \emph{\DUrole{n}{y}}, \emph{\DUrole{n}{line}\DUrole{o}{=}\DUrole{default_value}{False}}, \emph{\DUrole{n}{xloc}\DUrole{o}{=}\DUrole{default_value}{0.05}}, \emph{\DUrole{n}{yloc}\DUrole{o}{=}\DUrole{default_value}{0.1}}, \emph{\DUrole{n}{decPlace}\DUrole{o}{=}\DUrole{default_value}{2}}, \emph{\DUrole{n}{rSquared}\DUrole{o}{=}\DUrole{default_value}{False}}, \emph{\DUrole{n}{fontsize}\DUrole{o}{=}\DUrole{default_value}{15}}, \emph{\DUrole{n}{linecolor}\DUrole{o}{=}\DUrole{default_value}{\textquotesingle{}darkorange\textquotesingle{}}}, \emph{\DUrole{n}{lineDetails}\DUrole{o}{=}\DUrole{default_value}{False}}}{}
\pysigstopsignatures
\end{fulllineitems}

\index{rSquared() (xmmPipeline.plotAnal.plotAnal method)@\spxentry{rSquared()}\spxextra{xmmPipeline.plotAnal.plotAnal method}}

\begin{fulllineitems}
\phantomsection\label{\detokenize{xmmPipeline:xmmPipeline.plotAnal.plotAnal.rSquared}}
\pysigstartsignatures
\pysiglinewithargsret{\sphinxbfcode{\sphinxupquote{rSquared}}}{\emph{\DUrole{n}{y}}, \emph{\DUrole{n}{f}}}{}
\pysigstopsignatures
\sphinxAtStartPar
This can only be there between the y values and the fitted values f from the fitted line.

\end{fulllineitems}

\index{spearman() (xmmPipeline.plotAnal.plotAnal method)@\spxentry{spearman()}\spxextra{xmmPipeline.plotAnal.plotAnal method}}

\begin{fulllineitems}
\phantomsection\label{\detokenize{xmmPipeline:xmmPipeline.plotAnal.plotAnal.spearman}}
\pysigstartsignatures
\pysiglinewithargsret{\sphinxbfcode{\sphinxupquote{spearman}}}{\emph{\DUrole{n}{x}}, \emph{\DUrole{n}{y}}, \emph{\DUrole{n}{line}\DUrole{o}{=}\DUrole{default_value}{False}}, \emph{\DUrole{n}{xloc}\DUrole{o}{=}\DUrole{default_value}{0.05}}, \emph{\DUrole{n}{yloc}\DUrole{o}{=}\DUrole{default_value}{0.1}}, \emph{\DUrole{n}{decPlace}\DUrole{o}{=}\DUrole{default_value}{2}}, \emph{\DUrole{n}{rSquared}\DUrole{o}{=}\DUrole{default_value}{False}}, \emph{\DUrole{n}{fontsize}\DUrole{o}{=}\DUrole{default_value}{15}}, \emph{\DUrole{n}{linecolor}\DUrole{o}{=}\DUrole{default_value}{\textquotesingle{}darkorange\textquotesingle{}}}, \emph{\DUrole{n}{lineDetails}\DUrole{o}{=}\DUrole{default_value}{False}}}{}
\pysigstopsignatures
\end{fulllineitems}


\end{fulllineitems}



\subsection{xmmPipeline.pyXspec module}
\label{\detokenize{xmmPipeline:xmmpipeline-pyxspec-module}}

\subsection{xmmPipeline.rgsPipeline module}
\label{\detokenize{xmmPipeline:module-xmmPipeline.rgsPipeline}}\label{\detokenize{xmmPipeline:xmmpipeline-rgspipeline-module}}\index{module@\spxentry{module}!xmmPipeline.rgsPipeline@\spxentry{xmmPipeline.rgsPipeline}}\index{xmmPipeline.rgsPipeline@\spxentry{xmmPipeline.rgsPipeline}!module@\spxentry{module}}
\sphinxAtStartPar
01 April 2021:
—
Python routine to reduce the RGS data, obtain the Spectra and the light curve from it.
The SAS threads to be used, in order of execution are:
\begin{enumerate}
\sphinxsetlistlabels{\arabic}{enumi}{enumii}{}{.}%
\item {} 
\sphinxAtStartPar
\sphinxurl{https://www.cosmos.esa.int/web/xmm-newton/sas-thread-rgs}
Main thread to reduce the RGS data. Similar to ‘emproc’ and ‘epproc’.
Should do all the work and obtain the Spectra and the light curves.

\item {} 
\sphinxAtStartPar
\sphinxurl{https://www.cosmos.esa.int/web/xmm-newton/sas-thread-rgs2}
To correct the Source coordinates and the extraction masks, if the above
step was incorrect.

\item {} 
\sphinxAtStartPar
\sphinxurl{https://www.cosmos.esa.int/web/xmm-newton/sas-thread-timing}
For extraction of RGS light curves.

\end{enumerate}
\index{main() (in module xmmPipeline.rgsPipeline)@\spxentry{main()}\spxextra{in module xmmPipeline.rgsPipeline}}

\begin{fulllineitems}
\phantomsection\label{\detokenize{xmmPipeline:xmmPipeline.rgsPipeline.main}}
\pysigstartsignatures
\pysiglinewithargsret{\sphinxcode{\sphinxupquote{xmmPipeline.rgsPipeline.}}\sphinxbfcode{\sphinxupquote{main}}}{\emph{\DUrole{n}{args}}}{}
\pysigstopsignatures
\end{fulllineitems}

\index{printErrorMessage() (in module xmmPipeline.rgsPipeline)@\spxentry{printErrorMessage()}\spxextra{in module xmmPipeline.rgsPipeline}}

\begin{fulllineitems}
\phantomsection\label{\detokenize{xmmPipeline:xmmPipeline.rgsPipeline.printErrorMessage}}
\pysigstartsignatures
\pysiglinewithargsret{\sphinxcode{\sphinxupquote{xmmPipeline.rgsPipeline.}}\sphinxbfcode{\sphinxupquote{printErrorMessage}}}{\emph{\DUrole{n}{message}}}{}
\pysigstopsignatures
\end{fulllineitems}

\index{rgsObj (class in xmmPipeline.rgsPipeline)@\spxentry{rgsObj}\spxextra{class in xmmPipeline.rgsPipeline}}

\begin{fulllineitems}
\phantomsection\label{\detokenize{xmmPipeline:xmmPipeline.rgsPipeline.rgsObj}}
\pysigstartsignatures
\pysiglinewithargsret{\sphinxbfcode{\sphinxupquote{class\DUrole{w}{  }}}\sphinxcode{\sphinxupquote{xmmPipeline.rgsPipeline.}}\sphinxbfcode{\sphinxupquote{rgsObj}}}{\emph{\DUrole{n}{ra}}, \emph{\DUrole{n}{dec}}, \emph{\DUrole{n}{workdir}}, \emph{\DUrole{n}{sas\_dir}}, \emph{\DUrole{n}{headas}}, \emph{\DUrole{n}{sas\_ccfpath}}}{}
\pysigstopsignatures
\sphinxAtStartPar
Bases: \sphinxhref{https://docs.python.org/3/library/functions.html\#object}{\sphinxcode{\sphinxupquote{object}}}
\index{findObsIDs() (xmmPipeline.rgsPipeline.rgsObj method)@\spxentry{findObsIDs()}\spxextra{xmmPipeline.rgsPipeline.rgsObj method}}

\begin{fulllineitems}
\phantomsection\label{\detokenize{xmmPipeline:xmmPipeline.rgsPipeline.rgsObj.findObsIDs}}
\pysigstartsignatures
\pysiglinewithargsret{\sphinxbfcode{\sphinxupquote{findObsIDs}}}{}{}
\pysigstopsignatures\begin{quote}\begin{description}
\sphinxlineitem{Input}
\sphinxAtStartPar
Coordinates of the object and path to work directory.

\sphinxlineitem{Output}
\sphinxAtStartPar
List of obsIDs.

\end{description}\end{quote}

\end{fulllineitems}

\index{readPickleFile() (xmmPipeline.rgsPipeline.rgsObj method)@\spxentry{readPickleFile()}\spxextra{xmmPipeline.rgsPipeline.rgsObj method}}

\begin{fulllineitems}
\phantomsection\label{\detokenize{xmmPipeline:xmmPipeline.rgsPipeline.rgsObj.readPickleFile}}
\pysigstartsignatures
\pysiglinewithargsret{\sphinxbfcode{\sphinxupquote{readPickleFile}}}{}{}
\pysigstopsignatures
\sphinxAtStartPar
Reads the already obtained pickle file containing location parameters.

\end{fulllineitems}

\index{reduceRGSdata() (xmmPipeline.rgsPipeline.rgsObj method)@\spxentry{reduceRGSdata()}\spxextra{xmmPipeline.rgsPipeline.rgsObj method}}

\begin{fulllineitems}
\phantomsection\label{\detokenize{xmmPipeline:xmmPipeline.rgsPipeline.rgsObj.reduceRGSdata}}
\pysigstartsignatures
\pysiglinewithargsret{\sphinxbfcode{\sphinxupquote{reduceRGSdata}}}{}{}
\pysigstopsignatures
\sphinxAtStartPar
Reduces the RGS data using ‘rgsproc’ to obtain Spectra.
A folder names ‘RGS/’ is created in the work folder of each obsID directory.
See: \sphinxurl{https://www.cosmos.esa.int/web/xmm-newton/sas-thread-rgs}
\begin{quote}\begin{description}
\sphinxlineitem{Input}
\sphinxAtStartPar
None.

\sphinxlineitem{Output}
\sphinxAtStartPar
Source and Background RGS Spectra and the associated response matrices.

\end{description}\end{quote}

\end{fulllineitems}

\index{save\_results() (xmmPipeline.rgsPipeline.rgsObj method)@\spxentry{save\_results()}\spxextra{xmmPipeline.rgsPipeline.rgsObj method}}

\begin{fulllineitems}
\phantomsection\label{\detokenize{xmmPipeline:xmmPipeline.rgsPipeline.rgsObj.save_results}}
\pysigstartsignatures
\pysiglinewithargsret{\sphinxbfcode{\sphinxupquote{save\_results}}}{}{}
\pysigstopsignatures\begin{description}
\sphinxlineitem{Does two tasks,}\begin{description}
\sphinxlineitem{One, saves a pickle file at each obsID directory containing}
\sphinxAtStartPar
sourceCCDs, sourceLoc, backgroundLoc and otherSources for it.

\sphinxlineitem{Two, copies the Source, Background Events lists and other output files}
\sphinxAtStartPar
from each obsID directory to a results folder in the main directory.

\end{description}

\end{description}

\end{fulllineitems}


\end{fulllineitems}

\index{strToBool() (in module xmmPipeline.rgsPipeline)@\spxentry{strToBool()}\spxextra{in module xmmPipeline.rgsPipeline}}

\begin{fulllineitems}
\phantomsection\label{\detokenize{xmmPipeline:xmmPipeline.rgsPipeline.strToBool}}
\pysigstartsignatures
\pysiglinewithargsret{\sphinxcode{\sphinxupquote{xmmPipeline.rgsPipeline.}}\sphinxbfcode{\sphinxupquote{strToBool}}}{\emph{\DUrole{n}{s}}}{}
\pysigstopsignatures
\end{fulllineitems}



\subsection{Module contents}
\label{\detokenize{xmmPipeline:module-xmmPipeline}}\label{\detokenize{xmmPipeline:module-contents}}\index{module@\spxentry{module}!xmmPipeline@\spxentry{xmmPipeline}}\index{xmmPipeline@\spxentry{xmmPipeline}!module@\spxentry{module}}

\chapter{Indices and tables}
\label{\detokenize{index:indices-and-tables}}\begin{itemize}
\item {} 
\sphinxAtStartPar
\DUrole{xref,std,std-ref}{genindex}

\item {} 
\sphinxAtStartPar
\DUrole{xref,std,std-ref}{modindex}

\item {} 
\sphinxAtStartPar
\DUrole{xref,std,std-ref}{search}

\end{itemize}


\renewcommand{\indexname}{Python Module Index}
\begin{sphinxtheindex}
\let\bigletter\sphinxstyleindexlettergroup
\bigletter{x}
\item\relax\sphinxstyleindexentry{xmmPipeline}\sphinxstyleindexpageref{xmmPipeline:\detokenize{module-xmmPipeline}}
\item\relax\sphinxstyleindexentry{xmmPipeline.convert}\sphinxstyleindexpageref{xmmPipeline:\detokenize{module-xmmPipeline.convert}}
\item\relax\sphinxstyleindexentry{xmmPipeline.epicObj}\sphinxstyleindexpageref{xmmPipeline:\detokenize{module-xmmPipeline.epicObj}}
\item\relax\sphinxstyleindexentry{xmmPipeline.epicPileup}\sphinxstyleindexpageref{xmmPipeline:\detokenize{module-xmmPipeline.epicPileup}}
\item\relax\sphinxstyleindexentry{xmmPipeline.epicPipeline}\sphinxstyleindexpageref{xmmPipeline:\detokenize{module-xmmPipeline.epicPipeline}}
\item\relax\sphinxstyleindexentry{xmmPipeline.epicSpectra}\sphinxstyleindexpageref{xmmPipeline:\detokenize{module-xmmPipeline.epicSpectra}}
\item\relax\sphinxstyleindexentry{xmmPipeline.findOverlap}\sphinxstyleindexpageref{xmmPipeline:\detokenize{module-xmmPipeline.findOverlap}}
\item\relax\sphinxstyleindexentry{xmmPipeline.getSpectraFitsTable}\sphinxstyleindexpageref{xmmPipeline:\detokenize{module-xmmPipeline.getSpectraFitsTable}}
\item\relax\sphinxstyleindexentry{xmmPipeline.plotAnal}\sphinxstyleindexpageref{xmmPipeline:\detokenize{module-xmmPipeline.plotAnal}}
\item\relax\sphinxstyleindexentry{xmmPipeline.rgsPipeline}\sphinxstyleindexpageref{xmmPipeline:\detokenize{module-xmmPipeline.rgsPipeline}}
\end{sphinxtheindex}

\renewcommand{\indexname}{Index}
\printindex
\end{document}